\chapter{浮动基座建模与运动学}
\label{ch:quad_kin}

% \cnum（字体无关带圈数字）已上提到 common/preamble.tex，全书统一，勿在此重定义。

\begin{practice}[前置自测]
开始之前，先试答下列问题。若已能流畅作答，可只速读本章表格与速查卡；若卡住，正是本章要为你解决的。
\begin{enumerate}
  \item 四足机器人的电机只能直接控制 $12$ 个关节角，而我们想控制的却是“机身在世界中的位姿”与“足端落在哪里”。这两者之间靠什么数学对象连接？为什么说\emph{有运动学必先建坐标系}？
  \item 给定三个关节角求足端位置（正运动学）总有\emph{唯一}解且计算直接；反过来给定足端位置求关节角（逆运动学）却可能\emph{多解、无解、或奇异}。是什么造成了这种不对称？
  \item 足端位置 $\boldsymbol{p}$ 与关节角 $\boldsymbol{q}$ 之间是\emph{非线性}关系（不能写成矩阵乘 $\boldsymbol{p}=\mathbf{A}\boldsymbol{q}$）。可为什么它们的\emph{速度}之间却是线性的 $\dot{\boldsymbol{p}}=\mathbf{J}\dot{\boldsymbol{q}}$？这条“一阶微分把非线性变线性”的事实为什么如此有用？
  \item 同一个雅可比矩阵 $\mathbf{J}$，既出现在速度映射 $\dot{\boldsymbol{p}}=\mathbf{J}\dot{\boldsymbol{q}}$ 里，又出现在力映射 $\boldsymbol{\tau}=\mathbf{J}^\top\mathbf{F}$ 里。是什么物理原理让“速度”与“力”共用一个 $\mathbf{J}$？
  \item 一条 $7$ 自由度的臂去够一个 $6$ 维的目标位姿，会多出一个自由度。这“多出来的一维”在数学上落在 $\mathbf{J}$ 的什么子空间里？它对应机器人怎样的\emph{物理}运动？为什么说它是全身控制（WBC）多任务的本质？
  \item 当机器人接近\emph{奇异构型}时，$\mathbf{J}\mathbf{J}^\top$ 理论上仍可逆，可直接求逆 $\boldsymbol{q}=\mathbf{J}^{-1}\dot{\boldsymbol{p}}$ 却会让关节速度\emph{爆炸}。为什么？工程上怎么稳健化？
  \item 四足比定基机械臂多了一个会“在世界中自由漂移”的机身。如何用统一的运动学树框架，把这个浮动机身也当成一个普通关节来处理？
\end{enumerate}
\end{practice}

\section*{本章目标}
学完本章，你应当能够：
\begin{enumerate}
  \item \textbf{建立}四足机器人的完整坐标系体系（世界 / 机身 / 关节 / 连杆 / 足底系），理解“有运动学必先建坐标系”，并掌握旋转矩阵的三大用途与变换矩阵的两重作用（\cref{sec:qk-frames}）；
  \item \textbf{推导}旋转矩阵、ZYX 欧拉角及其逆解、以及角速度张量 $\dot{\mathbf{R}}\mathbf{R}^\top=\boldsymbol{\omega}^\wedge$ 的严格来历，并说清欧拉角为何只配人机交互、不宜长期描述姿态（\cref{sec:qk-frames}）；
  \item \textbf{用}虚拟 $6$ 自由度关节把浮动机身纳入运动学树，写出 $13$ 刚体的父节点集与浮动基座关节模型 $\mathbf{S}_1=\mathbf{I}_6$（\cref{sec:qk-floating}）；
  \item \textbf{计算}单腿正运动学（齐次变换连乘的闭式解）与单腿逆运动学（解耦三步：$\theta_1$ 投影 $+$ atan2、$\theta_3$ 余弦定理、$\theta_2$ 整理）（\cref{sec:qk-fk,sec:qk-ik}）；
  \item \textbf{构造}足端雅可比并理解“一阶微分把非线性变线性”这一全章方法论核心，进而由功率守恒证明运动学/静力学对偶 $\boldsymbol{\tau}=\mathbf{J}^\top\mathbf{F}$（\cref{sec:qk-jacobian}）；
  \item \textbf{掌握}Plücker 六维空间向量代数（运动 / 力向量、坐标变换 $\mathbf{X}/\mathbf{X}^*$、空间叉乘、空间惯量、运动方程）与递推正运动学、全身雅可比 $\mathbf{J}_1\dots\mathbf{J}_4$（\cref{sec:qk-plucker,sec:qk-wbjac}）；
  \item \textbf{证明}冗余系统的通解结构、零空间投影矩阵、Moore--Penrose 伪逆的最小范数性质（拉格朗日乘子法），并用阻尼最小二乘稳健化近奇异——把“奇异性与零空间”讲透作为 \cref{ch:quad_wbc}（WBC）的数学地基（\cref{sec:qk-nullspace}）；
  \item \textbf{整合}四足整机运动学：机身姿态$\leftrightarrow$足端位置、牵连速度 $+$ 关节速度 $\to$ 足端速度，并理解“机身系是瞬时重合的静止惯性系”这一认知关键（\cref{sec:qk-wholebody}）。
\end{enumerate}

\subsection*{本章知识导航}
本章是一条“\emph{先建坐标系，再用一阶微分把非线性运动学变成线性映射，最后用零空间理论为冗余全身控制铺地基}”的主线。结构如下：

\begin{center}
\begin{forest} navtreeLR
[{浮动基座建模\\与运动学}
  [{坐标系与位姿\\K1 坐标系即一切}
    [{旋转矩阵/欧拉角\\角速度张量}] ]
  [{浮动基座\\K10 虚拟 6-DoF 关节}
    [{运动学树\\13 刚体}] ]
  [{单腿 FK/IK\\K4 FK 易 IK 难}
    [{解耦三步逆解}] ]
  [{雅可比\\K5 一阶微分变线性}
    [{K6 静力学对偶 $\boldsymbol{\tau}{=}\mathbf{J}^\top\mathbf{F}$}] ]
  [{Plücker 空间向量\\递推 FK / 全身 $\mathbf{J}$}
    [{运动方程}] ]
  [{奇异与零空间\\K7/K8 WBC 地基}
    [{MP 伪逆证明}] ]
]
\end{forest}
\end{center}

\noindent\textbf{推荐路径}：\cref{sec:qk-frames} 建立坐标系与旋转表示，是一切的地基；\cref{sec:qk-floating} 用虚拟关节把浮动基座纳入统一框架，是四足区别于定基机械臂的唯一结构差别；\cref{sec:qk-fk,sec:qk-ik,sec:qk-jacobian} 是单腿运动学的“可手算”三件套（正解 / 逆解 / 雅可比），写控制器前必须吃透；\cref{sec:qk-plucker,sec:qk-wbjac} 用 Plücker 空间向量把单腿推广到整机递推与全身雅可比，是现代多刚体动力学的统一语言；\cref{sec:qk-nullspace} 是本章理论高潮——奇异性与零空间投影，直接喂给 \cref{ch:quad_wbc} 的多任务 WBC；\cref{sec:qk-wholebody} 把以上整合成整机运动学，衔接 \cref{ch:quad_est} 的足式状态估计。

\subsection*{前置知识桥接}
本章把前几部的几何工具直接用作地基：
\begin{itemize}
  \item \cref{ch:lie}（李群与李代数）给出旋转随时间演化的微分方程 $\dot{\mathbf{R}}=\mathbf{R}\,\boldsymbol{\omega}_b^\wedge$（机体系，右扰动主线）与 $\dot{\mathbf{R}}=\boldsymbol{\omega}^\wedge\mathbf{R}$（世界系），反对称算子 $(\cdot)^\wedge$、其逆 $(\cdot)^\vee$，以及矩阵指数 / 对数（轴角$\leftrightarrow$旋转矩阵）。本章角速度张量、单腿雅可比都建在这条微分方程上；
  \item \cref{ch:rigid_body}（刚体运动）给出旋转矩阵、$\SE(3)$ 齐次变换、ZYX 欧拉角与基本旋转 $\mathbf{R}_x,\mathbf{R}_y,\mathbf{R}_z$ 的定义，本章位姿运动学直接复用；
  \item \cref{ch:matderiv}（矩阵求导 / 雅可比）给出由向量函数 $\boldsymbol{f}(\boldsymbol{q})$ 对 $\boldsymbol{q}$ 求偏导得雅可比的机制，本章足端雅可比是它在单腿 FK 上的直接落地；
  \item \cref{ch:nlopt}（非线性优化 / QP）给出拉格朗日乘子法、最小二乘与 QP 的一般理论，本章 MP 伪逆证明、阻尼最小二乘是它的具体应用。
\end{itemize}

\subsection*{如果跳过本章会怎样}
\begin{itemize}
  \item 你在读 \cref{ch:quad_mpc}（单刚体 MPC）时会看到“MPC 算出足端力 $\mathbf{f}_i$、下游靠 $\boldsymbol{\tau}=\mathbf{J}^\top\mathbf{f}$ 转关节力矩”，却不知道这个 $\mathbf{J}$ 从哪来、为何力与速度共用它；
  \item 你在读 \cref{ch:quad_wbc}（WBC）时会看到“在零空间投影里做次级任务”“用阻尼伪逆求解优先级”，却不知道零空间 $\mathcal{N}(\mathbf{J})$、投影矩阵 $\mathbf{I}-\mathbf{J}^+\mathbf{J}$、MP 伪逆这些对象的几何意义与证明——只能照抄而不会调参、不会排错。
\end{itemize}

\begin{table}[H]\centering\small
\caption{本章阅读方式建议}
\begin{tabular}{lll}
\toprule
阅读方式 & 时间 & 适合谁 \\
\midrule
精读（含全部推导与练习） & 6--8 小时 & 首次系统学习、要做四足运动学 / WBC \\
速读（跳过 Plücker 递推堆叠） & 2.5 小时 & 有机械臂运动学基础、想对齐四足记号 \\
速查（只看小结的符号表 + 公式速查表） & 20 分钟 & 写运动学代码时回来查公式 / 维度 \\
\bottomrule
\end{tabular}
\end{table}

\begin{note}
\textbf{本章记号与约定.}\;
旋转矩阵记 ${}^{\mathcal{A}}\mathbf{R}_{\mathcal{B}}\in\SO(3)$（$\mathcal{A}$ 系看 $\mathcal{B}$ 系，列向量为 $\mathcal{B}$ 系坐标轴在 $\mathcal{A}$ 系的坐标），等价的左右下标式可写 $\mathbf{R}_{sb}$。
位姿姿态用 ZYX 欧拉角 $\boldsymbol{\Theta}=[\phi,\theta,\psi]^\top$（roll/pitch/yaw，\cref{ch:rigid_body}）。
反对称算子统一记 $(\cdot)^\wedge$（亦写 $[\cdot]_\times$）、其逆 $(\cdot)^\vee$（\cref{ch:lie}）。
单腿三关节角记 $\theta_1,\theta_2,\theta_3$（abad/hip/knee），它们即整机广义坐标 $\boldsymbol{q}$ 的对应腿关节分量 $q_i$；二者所指相同，单腿闭式推导用 $\theta_i$、整机递推框架（\cref{sec:qk-plucker} 起）用 $q_i$ 以贴合各自维度。
\textbf{旋转随时间演化以右扰动为主线} $\dot{\mathbf{R}}=\mathbf{R}\,\boldsymbol{\omega}_b^\wedge$（机体角速度），世界系空间角速度形式 $\dot{\mathbf{R}}=\boldsymbol{\omega}^\wedge\mathbf{R}$ 的换算见 \cref{note:qk-omega-frame}。
\textbf{六维空间向量（twist / wrench）采用角量在前约定} $[\boldsymbol{\omega}^\top,\boldsymbol{v}_O^\top]^\top$（\cref{sec:qk-plucker} 的 Plücker 空间向量代数标准分块）；全书 $\se(3)$ 位姿坐标则用平移在前 $\boldsymbol{\xi}=[\boldsymbol{\rho}^\top,\boldsymbol{\phi}^\top]^\top$，二者换算见 \cref{note:qk-twist-order,note:qk-plucker-order}。
浮动基座运动学树编号：虚拟固定基座 $=0$、机身 $=1$、$12$ 条腿连杆依次编号，父节点集见 \cref{eq:qk-lambda}。
\end{note}

% ════════════════════════════════════════════════════════════════
\section{坐标系与位姿描述：有运动学必先建坐标系}
\label{sec:qk-frames}

\begin{insight}[K1：坐标系即一切（全章开篇主线）]\label{ins:qk-frames-first}
全章第一句方法论是\textbf{“有运动学，必先建立坐标系”}。运动学研究的是“关节角$\to$足端$\to$机身”这条几何链上各量的相互换算，而\emph{任何一个几何量（位置、速度、力）都只有在某个坐标系里才有数值意义}。所以学四足运动学，第一件事不是写公式，而是把坐标系体系建清楚：四足按 \textbf{Body / Joint / Link} 分层——机身系、各关节系、各连杆系，再加世界系与足底系。机身系的 $xyz$ 轴约定为\emph{前 / 左 / 上}（$x$ 朝机器人前进方向、$z$ 朝上、$y$ 由右手系定为朝左）。一旦坐标系定清楚、记号统一，后面的旋转矩阵、雅可比、零空间都只是“在这些系之间搬运几何量”而已。本章所有看似复杂的上下标（如 ${}^{\mathcal{O}}\boldsymbol{p}_{i}$、${}^{\mathcal{B}}\mathbf{R}_{\mathcal{O}}$）都只是在回答一句话：\emph{这个量是哪个系里的、从哪个系看的}。
\end{insight}

\paragraph{动机：12 个电机是唯一可直接控的量。}
四足机器人有 $12$ 个驱动电机（每条腿三个：髋外摆 abad、髋俯仰 hip、膝 knee），这是\emph{唯一}可被直接施加指令的量。但我们真正关心的目标——“机身要走到哪里、姿态如何”“某只脚要落在地面哪一点”——都不是关节量，而是\emph{笛卡尔空间}的位姿\cite{unitree2023quadruped}。运动学正是连接这两端的桥：正运动学（FK）由关节角算足端 / 机身位姿；逆运动学（IK）反过来由期望位姿算关节角；一阶微分运动学把它们的速度联系起来；静力学把足端力换算为关节力矩。本章逐一打通这座桥。

\subsection{四足坐标系体系}
\label{subsec:qk-frame-system}

按 \cref{ins:qk-frames-first} 的分层，本章用到的坐标系汇总于 \cref{tab:qk-frames}。

\begin{table}[H]\centering\small
\caption{四足机器人坐标系一览}
\label{tab:qk-frames}
\begin{tabular}{p{0.13\textwidth} p{0.13\textwidth} p{0.62\textwidth}}
\toprule
记号 & 名称 & 定义与用途 \\
\midrule
$\mathcal{O}$（或 $\mathcal{W}$、$\{s\}$） & 世界 / 惯性系 & 固定于地面；机身位姿、足端绝对位置在此系表达 \\
$\mathcal{B}$（$\{b\}$） & 机身系 & 原点在机身（常取几何中心 / 重心），$xyz=$ 前 / 左 / 上；随机身运动（但见 \cref{ins:qk-static-body}）\\
$\{0\}\{1\}\{2\}$ & 单腿关节系 & 第 $i$ 腿三关节系：abad 绕 $x$、hip 绕 $y$、knee 绕 $y$ \\
$\{3\}$ & 足端系 & 足端（脚尖）所在；FK 的终点 \\
$\mathcal{Q}_i$ & 定向足底系 & 方向同世界、原点在第 $i$ 足底；全身雅可比的任务系 \\
\bottomrule
\end{tabular}
\end{table}

\begin{example}[四足机器人物理尺寸（贯穿全章数值例）]\label{exam:qk-dims}
取一组典型样机尺寸作为本章数值例的统一参数\cite{yu2021quadruped}：机身半长 $h_x=0.275\,\mathrm{m}$、机身半宽 $h_y=0.0625\,\mathrm{m}$，单腿三段长 $l_1=0.088\,\mathrm{m}$（abad 偏置）、$l_2=0.27\,\mathrm{m}$（大腿）、$l_3=0.27\,\mathrm{m}$（小腿）。A1 / Go1 的单腿则记为 $l_{abad},l_{hip},l_{knee}$\cite{unitree2023quadruped}。后文凡“代入数值”均默认此组参数。
\end{example}

\subsection{旋转矩阵：定义与三大用途}
\label{subsec:qk-rotation}

\paragraph{定义（坐标映射视角）。}
把 $\mathcal{B}$ 系的三个单位基向量在 $\mathcal{A}$ 系中的坐标列成列向量，即得旋转矩阵\cite{yu2021quadruped,unitree2023quadruped}：
\begin{equation}
  {}^{\mathcal{A}}\boldsymbol{r}_{AP}=\big[\,{}^{\mathcal{A}}\mathbf{e}_x^{\mathcal{B}}\ \ {}^{\mathcal{A}}\mathbf{e}_y^{\mathcal{B}}\ \ {}^{\mathcal{A}}\mathbf{e}_z^{\mathcal{B}}\,\big]\,{}^{\mathcal{B}}\boldsymbol{r}_{AP}
  ={}^{\mathcal{A}}\mathbf{R}_{\mathcal{B}}\,{}^{\mathcal{B}}\boldsymbol{r}_{AP}.
  \label{eq:qk-R-def}
\end{equation}
即“同一点 $P$ 在 $\mathcal{B}$ 系的坐标 ${}^{\mathcal{B}}\boldsymbol{r}_{AP}$，左乘 ${}^{\mathcal{A}}\mathbf{R}_{\mathcal{B}}$ 即换算成 $\mathcal{A}$ 系坐标”。等价的左右下标式 $\mathbf{R}_{sb}=[\hat{\boldsymbol{x}}_b\ \hat{\boldsymbol{y}}_b\ \hat{\boldsymbol{z}}_b]$ 直观体现“列即 $b$ 系坐标轴在 $s$ 系的坐标”，并给出极其好用的\textbf{下标消去原则}\cite{unitree2023quadruped}：
\begin{equation}
  {}^{\mathcal{A}}\mathbf{R}_{\mathcal{B}}\,{}^{\mathcal{B}}\mathbf{R}_{\mathcal{C}}={}^{\mathcal{A}}\mathbf{R}_{\mathcal{C}}.
  \label{eq:qk-R-cancel}
\end{equation}
相邻的内层下标 $\mathcal{B}$“消去”，恰似分数约分——记住这条，多坐标系链式换算就不会写错方向。

\paragraph{正交性与逆 $=$ 转置。}
旋转矩阵的列是单位正交基，故\cite{yu2021quadruped,unitree2023quadruped}
\begin{equation}
  {}^{\mathcal{A}}\mathbf{R}_{\mathcal{B}}^{\top}\,{}^{\mathcal{A}}\mathbf{R}_{\mathcal{B}}=\mathbf{I}_3,\qquad
  {}^{\mathcal{B}}\mathbf{R}_{\mathcal{A}}={}^{\mathcal{A}}\mathbf{R}_{\mathcal{B}}^{-1}={}^{\mathcal{A}}\mathbf{R}_{\mathcal{B}}^{\top},\qquad
  \det\mathbf{R}=+1\ (\text{右手系}).
  \label{eq:qk-R-orth}
\end{equation}

\begin{pitfall}[误区：对旋转矩阵求逆而非转置]
在计算机里对 $3\times3$ 矩阵做通用矩阵求逆既耗时又易引入数值误差，而旋转矩阵的逆\emph{恰好就是转置}\cite{unitree2023quadruped}。代码里凡需要 ${}^{\mathcal{B}}\mathbf{R}_{\mathcal{A}}$，一律写 \texttt{R.\allowbreak transpose()} 而非 \texttt{R.inverse()}。这是 \cref{eq:qk-R-orth} 在工程上的直接红线。
\end{pitfall}

\begin{note}[$\det\mathbf{R}=+1$ 的几何含义：旋转不缩放体积（连向后文操作度）]\label{note:qk-det-volume}
\cref{eq:qk-R-orth} 末项 $\det\mathbf{R}=+1$ 不是孤立的代数约定，而有直接几何含义\cite{unitree2023quadruped}：任一 $3\times3$ 矩阵 $\mathbf{A}$ 把单位立方体映成一个平行六面体，$\det\mathbf{A}$ 即\emph{该体的有向体积}（亦即 $\mathbf{A}$ 作为坐标变换时的\textbf{体积缩放比}）。例如 $\mathbf{A}=\tfrac13\mathbf{I}$ 把各棱缩到 $1/3$、体积缩到 $\det\mathbf{A}=1/27$。于是 $\det\mathbf{R}=+1$ 一句话说清两件事：\emph{(1)} 体积不变——\textbf{旋转既不缩放也不翻转物体}（正交矩阵本可 $\det=\pm1$，$+1$ 排除了含镜像的 $\mathrm{O}(3)$ 反射，故 $\mathbf{R}\in\SO(3)$ 是“纯转动”）；\emph{(2)} 正号保证右手系不被翻成左手系。这条“行列式 $=$ 体积比”的直觉在本章后文反复复用：奇异性度量\emph{操作度} $w=\sqrt{\det(\mathbf{J}\mathbf{J}^\top)}=|\det\mathbf{J}|$（\cref{exam:qk-singular}）正是把 $\mathbf{J}$ 看作“关节速度立方体$\to$足端速度平行六面体”的体积缩放比——$w\to0$ 即该六面体被压扁（某方向足端速度无法生成），就是奇异。
\end{note}

\paragraph{基本旋转矩阵。}
绕 $x,y,z$ 轴转 $\varphi$ 的基本旋转为（$c_\varphi=\cos\varphi,\ s_\varphi=\sin\varphi$）\cite{yu2021quadruped,unitree2023quadruped}：
\begin{equation}
  \mathbf{R}_x(\varphi)=\begin{bmatrix}1&0&0\\0&c_\varphi&-s_\varphi\\0&s_\varphi&c_\varphi\end{bmatrix},\
  \mathbf{R}_y(\varphi)=\begin{bmatrix}c_\varphi&0&s_\varphi\\0&1&0\\-s_\varphi&0&c_\varphi\end{bmatrix},\
  \mathbf{R}_z(\varphi)=\begin{bmatrix}c_\varphi&-s_\varphi&0\\s_\varphi&c_\varphi&0\\0&0&1\end{bmatrix}.
  \label{eq:qk-basic-R}
\end{equation}

\begin{insight}[K2：旋转矩阵三用途 + 变换矩阵两重作用（认知锚点）]\label{ins:qk-three-uses}
一个统一记号的认知锚点是：\textbf{旋转矩阵有且仅有三大用途}——(1) \emph{描述姿态}（一个系相对另一个系的朝向，即 ${}^{\mathcal{A}}\mathbf{R}_{\mathcal{B}}$ 本身就是 $\mathcal{B}$ 的姿态）；(2) \emph{坐标变换}（同一向量在两系坐标间换算，\cref{eq:qk-R-def}）；(3) \emph{向量旋转}（把一个向量在\emph{同一系}内主动转一个角度，$\boldsymbol{r}'=\mathbf{R}\boldsymbol{r}$）。这三者数学上都是“左乘一个 $\mathbf{R}$”，却是三种完全不同的物理语义，混淆它们是初学者最大的坑。

进一步，齐次变换矩阵 $\mathbf{T}$ 也有\emph{两重作用}：作\emph{坐标系映射}（把一个系的位姿表达到另一个系），或作\emph{对向量的操作}（先转后移 / 先移后转地搬动一个点）。关键差异在于：\textbf{对向量操作时，平移与旋转的先后顺序影响结果}（不可交换）；而对坐标系做变换时则\emph{无此先后问题}（结果由两系的相对位姿唯一确定）。记住这条，连乘 $\mathbf{T}$ 时就不会把“描述位姿”与“移动物体”两种意图搞混。
\end{insight}

\paragraph{连续旋转与“固定轴 / 动轴”左右乘法则。}
连续旋转满足结合律、不满足交换律\cite{yu2021quadruped,unitree2023quadruped}：
\begin{equation}
  {}^{\mathcal{A}}\mathbf{R}_{\mathcal{C}}={}^{\mathcal{A}}\mathbf{R}_{\mathcal{B}}\,{}^{\mathcal{B}}\mathbf{R}_{\mathcal{C}}.
  \label{eq:qk-R-compose}
\end{equation}
一条好记的法则是\cite{unitree2023quadruped}：\emph{绕固定（世界）轴}转动 $\Rightarrow$ 新矩阵\textbf{左乘}；\emph{绕动（自身）轴}转动 $\Rightarrow$ 新矩阵\textbf{右乘}。ZYX 欧拉角的 $\mathbf{R}_z(\psi)\mathbf{R}_y(\theta)\mathbf{R}_x(\phi)$（见下）正是“绕固定世界轴依次 $z,y,x$ 左乘”的产物。

\paragraph{轴角与罗德里格斯公式。}
绕单位轴 $\hat{\boldsymbol{u}}$ 逆时针转 $\theta$ 的旋转矩阵由罗德里格斯公式给出\cite{yu2021quadruped}：
\begin{equation}
  \mathbf{R}=\mathbf{I}+(\sin\theta)\,\hat{\boldsymbol{u}}^{\wedge}+(1-\cos\theta)\,(\hat{\boldsymbol{u}}^{\wedge})^2,
  \qquad
  \hat{\boldsymbol{u}}^{\wedge}=\begin{bmatrix}0&-u_z&u_y\\u_z&0&-u_x\\-u_y&u_x&0\end{bmatrix}.
  \label{eq:qk-rodrigues}
\end{equation}
其中 $(\hat{\boldsymbol{u}}^{\wedge})^2$ 是\emph{反对称矩阵的平方}（一个矩阵自乘），\emph{不是}对反对称矩阵的每个元素逐个平方——后者会得到一个非反对称、无几何意义的矩阵，须留意区分。该式等价地写成\textbf{指数坐标} $\mathbf{R}=\exp(\theta\hat{\boldsymbol{u}}^{\wedge})$，配矩阵对数作逆映射\cite{unitree2023quadruped}；这套指数 / 对数映射的完整理论见 \cref{ch:lie}，此处给闭式以备单腿推导之用。

\subsection{ZYX 欧拉角及其逆解}
\label{subsec:qk-euler}

\paragraph{欧拉角 $\to$ 旋转矩阵。}
取 ZYX 序（先绕 $z$ 偏航 $\psi$、再绕 $y$ 俯仰 $\theta$、最后绕 $x$ 横滚 $\phi$，均绕固定世界轴故左乘）\cite{yu2021quadruped}：
\begin{equation}
  \boldsymbol{\Theta}=[\phi\ \theta\ \psi]^\top,\qquad
  \mathbf{R}(\boldsymbol{\Theta})=\mathbf{R}_z(\psi)\mathbf{R}_y(\theta)\mathbf{R}_x(\phi).
  \label{eq:qk-euler-to-R}
\end{equation}
展开为（$c_\bullet=\cos\bullet,s_\bullet=\sin\bullet$）\cite{yu2021quadruped}：
\begin{equation}
  \mathbf{R}=\begin{bmatrix}
    c_\theta c_\psi & c_\psi s_\phi s_\theta-c_\phi s_\psi & s_\phi s_\psi+c_\phi c_\psi s_\theta\\
    c_\theta s_\psi & c_\phi c_\psi+s_\phi s_\theta s_\psi & c_\phi s_\theta s_\psi-c_\psi s_\phi\\
    -s_\theta & c_\theta s_\phi & c_\phi c_\theta
  \end{bmatrix}.
  \label{eq:qk-euler-R-expand}
\end{equation}

\paragraph{旋转矩阵 $\to$ 欧拉角（逆解）。}
记 $\mathbf{R}=[c_{ij}]$，由 \cref{eq:qk-euler-R-expand} 比对元素：$c_{31}=-s_\theta$ 给 $\theta$，第三列与第一列比值给 $\phi,\psi$\cite{yu2021quadruped}：
\begin{equation}
  \boldsymbol{\Theta}=\begin{bmatrix}\phi\\\theta\\\psi\end{bmatrix}
  =\begin{bmatrix}
    \operatorname{atan2}(c_{32},\,c_{33})\\[2pt]
    \operatorname{atan2}\!\big(-c_{31},\,\sqrt{c_{32}^2+c_{33}^2}\big)\\[2pt]
    \operatorname{atan2}(c_{21},\,c_{11})
  \end{bmatrix}.
  \label{eq:qk-R-to-euler}
\end{equation}

\begin{note}[逆解必须用 atan2 而非单参数 atan]\label{note:qk-ocr-atan2}
\cref{eq:qk-R-to-euler} 的三个角分量一律用双参数 $\operatorname{atan2}(y,x)$ 而非单参数 $\operatorname{atan}$。单参数反正切的值域只有 $(-\pi/2,\pi/2)$、无法区分象限，会丢失姿态的符号信息导致逆解错误；而 $\operatorname{atan2}$ 值域 $(-\pi,\pi]$、依分子分母各自的符号唯一定出象限。同样的理由贯穿后文单腿逆解（\cref{eq:qk-ik-theta1,eq:qk-ik-theta2}）。
\end{note}

\begin{insight}[K3：欧拉角只配人机交互，不宜长期描述姿态（反面教材）]\label{ins:qk-euler-only-hmi}
欧拉角宜被定位为\emph{反面教材}：\textbf{欧拉角直观（人能一眼读出 roll/pitch/yaw），故适合人机交互（遥控器、参数面板、可视化），但不适用于机器人姿态的\emph{长期描述}}。原因有二：(1) \textbf{万向锁奇异}——当俯仰 $\theta\to\pm\pi/2$（$c_\theta\to0$），\cref{eq:qk-R-to-euler} 的 $\sqrt{c_{32}^2+c_{33}^2}\to0$，$\phi,\psi$ 不再可唯一确定，丢失一个自由度；(2) \textbf{无法平滑插值 / 迭代}——三个角分量的旋转轴随姿态变化且彼此不正交，直接对欧拉角插值会得到扭曲的中间姿态。因此姿态的长期表达与积分应当用\emph{旋转矩阵或四元数}（本书全局采用 Hamilton 四元数主线），欧拉角只在输入 / 显示这一层出现。这条结论在 \cref{ch:quad_mpc} 也复现：单刚体 MPC 用欧拉角换得线性时变模型，代价就是大俯仰下的奇异，须改用几何 MPC。
\end{insight}

\subsection{齐次变换矩阵}
\label{subsec:qk-se3}

把旋转与平移合并为 $\SE(3)$ 齐次变换\cite{unitree2023quadruped}：
\begin{equation}
  \mathbf{T}=\begin{bmatrix}\mathbf{R}&\boldsymbol{t}\\\mathbf{0}^\top&1\end{bmatrix}\in\SE(3),
  \qquad
  \mathbf{T}^{-1}=\begin{bmatrix}\mathbf{R}^\top&-\mathbf{R}^\top\boldsymbol{t}\\\mathbf{0}^\top&1\end{bmatrix}.
  \label{eq:qk-T}
\end{equation}
向量须齐次扩展（末位补 $1$）方可被 $\mathbf{T}$ 作用，$\mathbf{T}$ 可像旋转矩阵一样按 \cref{eq:qk-R-compose} 连乘——这是单腿正运动学的核心机械操作（\cref{sec:qk-fk}）。逆 $\mathbf{T}^{-1}$ 的闭式（\cref{eq:qk-T} 右）同样避免了通用矩阵求逆。

\begin{note}[se(3) 平移在前的换算]\label{note:qk-twist-order}
\cref{eq:qk-T} 中 $\mathbf{T}$ 本身的块结构与 twist 排序约定无关。本书的 $\SE(3)$ 李代数\emph{位姿坐标}（twist）$\boldsymbol{\xi}$ 统一采用 \textbf{平移在前} $\boldsymbol{\xi}=[\boldsymbol{\rho}^\top,\boldsymbol{\phi}^\top]^\top$（线量在前、角量在后）；而 \cref{sec:qk-plucker} 的 Plücker 空间向量（twist/wrench）采用 \textbf{角量在前} $[\boldsymbol{\omega}^\top,\boldsymbol{v}_O^\top]^\top$（空间向量代数的标准分块）。二者仅差上下 $3$ 维分块的交换，对应的伴随 / Plücker 变换矩阵 $\mathbf{X}$ 也按分块转置对应交换，详见 \cref{note:qk-plucker-order}。
\end{note}

\subsection{角速度张量：\texorpdfstring{$\dot{\mathbf{R}}\mathbf{R}^\top=\boldsymbol{\omega}^\wedge$}{Rdot R^T = omega^} 的严格推导}
\label{subsec:qk-angular-vel}

本小节给出“角速度张量”的严格推导\cite{yu2021quadruped}——这是把姿态 $\mathbf{R}$ 与角速度 $\boldsymbol{\omega}$ 联系起来、后续雅可比与整机速度公式的地基。

\paragraph{反对称算子。}
先定义反对称算子（叉乘的矩阵形式）\cite{yu2021quadruped,unitree2023quadruped}：
\begin{equation}
  \boldsymbol{a}^{\wedge}=[a_1\ a_2\ a_3]^\top{}^{\wedge}
  =\begin{bmatrix}0&-a_3&a_2\\a_3&0&-a_1\\-a_2&a_1&0\end{bmatrix},
  \qquad \boldsymbol{a}\times\boldsymbol{b}=\boldsymbol{a}^{\wedge}\boldsymbol{b}.
  \label{eq:qk-skew}
\end{equation}

\paragraph{角速度张量定义与反对称性。}
考虑刚接于刚体 $\mathcal{B}$ 上的点 $P$（其在 $\mathcal{B}$ 系坐标 ${}^{\mathcal{B}}\boldsymbol{r}_{AP}$ 为常数），对 \cref{eq:qk-R-def} 两边求时间导数\cite{yu2021quadruped}：
\begin{equation}
  {}^{\mathcal{A}}\dot{\boldsymbol{r}}_{AP}
  ={}^{\mathcal{A}}\dot{\mathbf{R}}_{\mathcal{B}}\,{}^{\mathcal{B}}\boldsymbol{r}_{AP}
  ={}^{\mathcal{A}}\dot{\mathbf{R}}_{\mathcal{B}}\,{}^{\mathcal{A}}\mathbf{R}_{\mathcal{B}}^\top\,{}^{\mathcal{A}}\boldsymbol{r}_{AP}
  =\underbrace{\big({}^{\mathcal{A}}\dot{\mathbf{R}}_{\mathcal{B}}\,{}^{\mathcal{A}}\mathbf{R}_{\mathcal{B}}^\top\big)}_{{}^{\mathcal{A}}\boldsymbol{\Omega}_{\mathcal{B}}}\,{}^{\mathcal{A}}\boldsymbol{r}_{AP},
  \label{eq:qk-angvel-1}
\end{equation}
其中代入了 ${}^{\mathcal{B}}\boldsymbol{r}_{AP}={}^{\mathcal{A}}\mathbf{R}_{\mathcal{B}}^\top\,{}^{\mathcal{A}}\boldsymbol{r}_{AP}$（由 \cref{eq:qk-R-orth}）。称 ${}^{\mathcal{A}}\boldsymbol{\Omega}_{\mathcal{B}}={}^{\mathcal{A}}\dot{\mathbf{R}}_{\mathcal{B}}\,{}^{\mathcal{A}}\mathbf{R}_{\mathcal{B}}^\top$ 为\textbf{角速度张量}。对正交性恒等式 $\mathbf{R}\mathbf{R}^\top=\mathbf{I}$ 求导：
\begin{equation}
  \frac{\mathrm{d}}{\mathrm{d}t}(\mathbf{R}\mathbf{R}^\top)
  =\dot{\mathbf{R}}\mathbf{R}^\top+\mathbf{R}\dot{\mathbf{R}}^\top
  =\boldsymbol{\Omega}+\boldsymbol{\Omega}^\top
  =\mathbf{0}
  \;\Longrightarrow\;
  \boldsymbol{\Omega}=-\boldsymbol{\Omega}^\top.
  \label{eq:qk-angvel-skew}
\end{equation}
故角速度张量\emph{反对称}，必可写成某向量 $\boldsymbol{\omega}$ 的反对称算子\cite{yu2021quadruped}：
\begin{equation}
  {}^{\mathcal{A}}\boldsymbol{\Omega}_{\mathcal{B}}
  =\begin{bmatrix}0&-\omega_z&\omega_y\\\omega_z&0&-\omega_x\\-\omega_y&\omega_x&0\end{bmatrix}
  ={}^{\mathcal{A}}\boldsymbol{\omega}_{\mathcal{AB}}^{\wedge},
  \qquad\Longrightarrow\qquad
  {}^{\mathcal{A}}\dot{\mathbf{R}}_{\mathcal{B}}={}^{\mathcal{A}}\boldsymbol{\omega}_{\mathcal{AB}}^{\wedge}\,{}^{\mathcal{A}}\mathbf{R}_{\mathcal{B}}.
  \label{eq:qk-Rdot-world}
\end{equation}
代回 \cref{eq:qk-angvel-1} 即得刚体上点的速度场 ${}^{\mathcal{A}}\dot{\boldsymbol{r}}_{AP}={}^{\mathcal{A}}\boldsymbol{\omega}_{\mathcal{AB}}\times{}^{\mathcal{A}}\boldsymbol{r}_{AP}$——纯转动贡献的线速度。

\begin{note}[世界系 / 机体系角速度与右扰动主线的换算]\label{note:qk-omega-frame}
\cref{eq:qk-Rdot-world} 是\emph{世界（左）形式} $\dot{\mathbf{R}}=\boldsymbol{\omega}_w^{\wedge}\mathbf{R}$，其中 $\boldsymbol{\omega}_w={}^{\mathcal{A}}\boldsymbol{\omega}_{\mathcal{AB}}$ 是世界系下表达的空间角速度。本书几何主线（\cref{ch:lie}）以\emph{右扰动}为主，对应\emph{机体角速度}形式 $\dot{\mathbf{R}}=\mathbf{R}\,\boldsymbol{\omega}_b^{\wedge}$，其中 ${}^{\mathcal{B}}\boldsymbol{\omega}_{\mathcal{AB}}^{\wedge}={}^{\mathcal{B}}\mathbf{R}_{\mathcal{A}}\,{}^{\mathcal{A}}\dot{\mathbf{R}}_{\mathcal{B}}$\cite{yu2021quadruped}。二者换算为
\[
  \boldsymbol{\omega}_w=\mathbf{R}\,\boldsymbol{\omega}_b,
  \qquad(\text{由 }\boldsymbol{\omega}_w^{\wedge}=\dot{\mathbf{R}}\mathbf{R}^\top,\ \boldsymbol{\omega}_b^{\wedge}=\mathbf{R}^\top\dot{\mathbf{R}},\ \text{及 }\mathbf{R}\boldsymbol{a}^{\wedge}\mathbf{R}^\top=(\mathbf{R}\boldsymbol{a})^{\wedge}).
\]
本章在整机速度（\cref{sec:qk-wholebody}）多用机体角速度（右扰动主线），在角速度张量推导这一段保留世界系形式以贴合该推导的自然写法；读者只需记住二者差一个 $\mathbf{R}$。
\end{note}

\begin{pitfall}[误区：把角速度当成欧拉角的导数 $\boldsymbol{\omega}=\dot{\boldsymbol{\Theta}}$]
不能把角速度简单等同于欧拉角对时间的导数。欧拉角三分量的旋转轴随姿态变化、彼此\emph{不正交}，以它们为基描述角速度违背物理。角速度由 \cref{eq:qk-Rdot-world} 的 $\dot{\mathbf{R}}\mathbf{R}^\top$ 严格定义，它与欧拉角率 $\dot{\boldsymbol{\Theta}}$ 之间差一个\emph{依赖姿态}的映射矩阵（其完整形式与小角近似见 \cref{ch:quad_mpc} 的凸化推导）。混淆二者是姿态控制最常见的错误来源。
\end{pitfall}

% ════════════════════════════════════════════════════════════════
\section{浮动基座与运动学树}
\label{sec:qk-floating}

\paragraph{动机：四足与定基机械臂的唯一结构差别。}
定基机械臂的底座固定在世界中，运动学链从一个已知不动的根出发。四足机器人却多了一个会“在世界中自由漂移”的机身——它有 $6$ 个不受电机直接驱动的自由度（$3$ 平移 $+$ $3$ 转动）。这就是\emph{浮动基座}（floating base）。如何把这个“漂浮的根”纳入与机械臂统一的运动学框架？答案是一个优雅的技巧。

\begin{insight}[K10：浮动基座 $=$ 虚拟 6-DoF 关节（统一记号的关键）]\label{ins:qk-floating-virtual}
处理浮动基座的统一手法是：在\emph{虚拟}的固定基座（编号 $0$，固定于世界、永不运动）与真实机身（编号 $1$）之间，插入一个\textbf{虚拟的 $6$ 自由度关节}\cite{yu2021quadruped}。于是机身的位姿就成了这个虚拟关节的“关节变量”：
\begin{equation}
  \boldsymbol{q}_1=\big[\,\boldsymbol{\Theta}^\top\ \ {}^{\mathcal{O}}\boldsymbol{p}_{com}^\top\,\big]^\top\in\mathbb{R}^6,
  \qquad {}^{1}\mathbf{S}_1=\mathbf{I}_{6\times6}.
  \label{eq:qk-floating-joint}
\end{equation}
其关节运动子空间 $\mathbf{S}_1=\mathbf{I}_6$（六个自由度全开放，无任何约束）。\textbf{这是浮动基座与定基机械臂\emph{唯一}的结构差别}：机械臂的根关节是固定的（$\mathbf{S}_0$ 为空），四足的根关节是一个全开放的 $6$-DoF 关节。一旦这样处理，整机就退化为一棵普通的运动学树，所有递推算法（\cref{sec:qk-plucker}）对浮动基座与定基系统\emph{完全一致}——这正是统一记号、复用机械臂理论的钥匙。
\end{insight}

\paragraph{运动学树与父节点集。}
四足整机建模为 $N=13$ 个刚体（$1$ 机身 $+$ $4$ 腿 $\times$ $3$ 连杆）。每个刚体 $i$ 有唯一父节点 $\lambda(i)$，构成\emph{运动学树}\cite{yu2021quadruped}：
\begin{equation}
  \lambda=\{\lambda(1),\dots,\lambda(N)\}=\{0,\,1,2,3,\,1,5,6,\,1,8,9,\,1,11,12\},\qquad N=13.
  \label{eq:qk-lambda}
\end{equation}
读法：刚体 $1$（机身）的父节点是 $0$（虚拟固定基座）；刚体 $2,3$ 的父节点链是 $1\to2\to3$（第一条腿三连杆）；刚体 $4$ 的父节点是 $1$（第二条腿首连杆挂在机身上），依此类推。四条腿首连杆的父节点都是 $1$（机身），体现“机身是所有腿的共同根”。

\begin{figure}[ht]
\centering
\begin{tikzpicture}[>=Stealth, node distance=6mm,
  bd/.style={draw, circle, fill=main!12, draw=main!60!black, minimum size=8mm, font=\scriptsize},
  vb/.style={draw, circle, fill=gray!15, draw=gray!60!black, minimum size=8mm, font=\scriptsize},
  lk/.style={draw, circle, fill=second!10, draw=second!60!black, minimum size=7mm, font=\scriptsize}]
  \node[vb] (b0) at (0,3) {$0$};
  \node[bd] (b1) at (0,1.8) {$1$};
  \draw[->, thick, gray] (b0)--(b1) node[midway,right,font=\tiny]{虚拟 6-DoF $\mathbf{S}_1{=}\mathbf{I}_6$};
  % four legs
  \foreach \base/\a/\bx in {2/-4.5/RF, 5/-1.5/LF, 8/1.5/RH, 11/4.5/LH}{
    \pgfmathtruncatemacro{\mid}{\base+1}
    \pgfmathtruncatemacro{\foot}{\base+2}
    \node[lk] (l\base) at (\a,0.4) {$\base$};
    \node[lk] (l\mid) at (\a,-0.8) {$\mid$};
    \node[lk] (l\foot) at (\a,-2.0) {$\foot$};
    \draw[->, thick, second!60!black] (b1)--(l\base);
    \draw[->, thick, second!60!black] (l\base)--(l\mid);
    \draw[->, thick, second!60!black] (l\mid)--(l\foot);
    \node[font=\tiny] at (\a,-2.5) {\bx};
  }
\end{tikzpicture}
\caption{四足运动学树：虚拟固定基座 $0$（灰）经虚拟 $6$ 自由度关节连机身 $1$（蓝），机身再分出四条腿（RF/LF/RH/LH 各三连杆，编号见 \cref{eq:qk-lambda}）。足端刚体为 $\{4,7,10,13\}$。浮动基座被统一为一个 $\mathbf{S}_1=\mathbf{I}_6$ 的普通树关节。}
\label{fig:qk-tree}
\end{figure}

\begin{note}[历史脉络：从 DH 到旋量到 Featherstone 空间向量]\label{note:qk-modeling}
运动学建模范式经历了三代演进，理解这条历史线有助于定位本章的工具选择：
\begin{itemize}
  \item \textbf{DH 参数法（Denavit--Hartenberg）}：机械臂运动学的经典，用 $4$ 个参数（连杆长 / 扭角 / 偏置 / 关节角）描述相邻连杆系。优点是参数最少，缺点是坐标系放置规则繁琐、对树状 / 浮动基座不自然，今多作历史背景\cite{yu2021quadruped}。
  \item \textbf{齐次变换连乘（旋量 / 指数坐标）}：直接用 $\mathbf{T}_{i-1,i}$ 连乘，几何意义清晰、工程实用，单腿 FK（\cref{sec:qk-fk}）即用此法\cite{unitree2023quadruped}。
  \item \textbf{Featherstone 空间向量（Plücker 代数）}：用六维空间向量统一描述运动与力，配运动学树递推，是现代多刚体动力学的统一框架，本章整机递推（\cref{sec:qk-plucker}）采此法\cite{yu2021quadruped}。
\end{itemize}
四足主用后两者：齐次变换给单腿“可手算”的闭式，Plücker 给整机“可扩展到 $N$ 刚体”的递推，二者互验。
\end{note}

% ════════════════════════════════════════════════════════════════
\section{单腿正运动学}
\label{sec:qk-fk}

\begin{insight}[K4：FK 易、IK 难（运动学的根本不对称）]\label{ins:qk-fk-easy-ik-hard}
正逆运动学的本质区别在于：\textbf{正运动学（FK）由因推果——给定关节角求足端位置，是唯一解、计算直接}；\textbf{逆运动学（IK）由果推因——给定足端位置求关节角，可能多解、无解、或奇异解}。FK 只是把关节角代入确定的几何函数 $\boldsymbol{f}(\boldsymbol{q})$；IK 却要解一组非线性方程，解的存在性与个数取决于目标点是否落在\emph{可达工作空间}（Reachable workspace，能以\emph{某}姿态够到的点集）与\emph{灵巧工作空间}（Dexterous workspace，能以\emph{任意}姿态够到的点集）内。所以做 IK 之前应先判断目标点的可达性，这是先验地避免“无解”的方法。
\end{insight}

\paragraph{二维引子。}
先看平面二连杆（连杆长 $l_1,l_2$，关节角 $\theta_1,\theta_2$）的足端位置\cite{unitree2023quadruped}：
\begin{equation}
  x_P=l_1\cos\theta_1+l_2\cos(\theta_1+\theta_2),
  \qquad
  y_P=l_1\sin\theta_1+l_2\sin(\theta_1+\theta_2).
  \label{eq:qk-2d-fk}
\end{equation}
注意 $x_P,y_P$ 是 $\theta_1,\theta_2$ 的\emph{非线性}函数（含 $\cos,\sin$ 的复合）——这一非线性正是 \cref{ins:qk-fk-easy-ik-hard} 中“IK 难”的根源，也是 \cref{sec:qk-jacobian} 要用一阶微分“线性化”的对象。

\begin{derivation}[二维引子的齐次变换法：为何要换用 $\mathbf{T}$ 连乘]\label{deriv:qk-2d-T}
\cref{eq:qk-2d-fk} 是\emph{直接看图}用三角几何凑出的，平面下尚可，但三维下连杆方向纠缠、几乎无法直接“看”出。源书\cite{unitree2023quadruped}于此处先在二维上验证另一条可\emph{机械推广}到三维的路线——齐次变换连乘——再放心地把它用于 \cref{eq:qk-3d-fk}。建三系：基座 $\{0\}$、随连杆 $1$ 转的 $\{1\}$（与 $\{0\}$ 原点重合、夹角 $\theta_1$）、随连杆 $2$ 的 $\{2\}$（相对 $\{1\}$ 夹角 $\theta_2$、原点沿连杆 $1$ 平移 $l_1$）。相邻齐次变换（二维 $\SE(2)$，平移列即下一原点在本系坐标）为
\begin{equation}
  \mathbf{T}_{01}=\begin{bmatrix}c_1&-s_1&0\\ s_1&c_1&0\\ 0&0&1\end{bmatrix},\qquad
  \mathbf{T}_{12}=\begin{bmatrix}c_2&-s_2&l_1\\ s_2&c_2&0\\ 0&0&1\end{bmatrix}.
  \label{eq:qk-2d-T0112}
\end{equation}
按下标消去 \cref{eq:qk-R-cancel} 连乘，旋转块用和角公式 $c_1c_2-s_1s_2=c_{12}$、$s_1c_2+c_1s_2=s_{12}$ 合并：
\begin{equation}
  \mathbf{T}_{02}=\mathbf{T}_{01}\mathbf{T}_{12}
  =\begin{bmatrix}c_{12}&-s_{12}&l_1 c_1\\ s_{12}&c_{12}&l_1 s_1\\ 0&0&1\end{bmatrix}.
  \label{eq:qk-2d-T02}
\end{equation}
足端 $P$ 在 $\{2\}$ 系的齐次坐标是常量 $[\,l_2,\,0,\,1\,]^\top$，左乘 $\mathbf{T}_{02}$ 抬到 $\{0\}$ 系：
\begin{equation}
  \begin{bmatrix}x_P\\ y_P\\ 1\end{bmatrix}
  =\mathbf{T}_{02}\begin{bmatrix}l_2\\ 0\\ 1\end{bmatrix}
  =\begin{bmatrix}l_2 c_{12}+l_1 c_1\\ l_2 s_{12}+l_1 s_1\\ 1\end{bmatrix},
  \label{eq:qk-2d-fk-T}
\end{equation}
与直接几何的 \cref{eq:qk-2d-fk} \emph{逐项一致}。\textbf{要点}：齐次变换法没有用到任何“看图凑”的技巧，纯靠“常量足端坐标 $\to$ 逐级左乘相邻 $\mathbf{T}$”这一\emph{机械}操作，故能原样搬到三维（\cref{eq:qk-fk-step1,eq:qk-fk-step2,eq:qk-fk-step3}）。这正是下文舍几何、取 $\mathbf{T}$ 连乘的理由。
\end{derivation}

\paragraph{三维单腿齐次变换。}
四足单腿三关节：机身关节（abad）绕 $x$、大腿关节（hip）绕 $y$、小腿关节（knee）绕 $y$。取坐标系 $\{0\}\{1\}\{2\}$ 原点重合（简化），相邻系齐次变换为\cite{unitree2023quadruped}：
\begin{equation}
  \mathbf{T}_{01}=\begin{bmatrix}1&0&0&0\\0&c_1&-s_1&0\\0&s_1&c_1&0\\0&0&0&1\end{bmatrix},\
  \mathbf{T}_{12}=\begin{bmatrix}c_2&0&s_2&0\\0&1&0&0\\-s_2&0&c_2&0\\0&0&0&1\end{bmatrix},\
  \mathbf{T}_{23}=\begin{bmatrix}c_3&0&s_3&0\\0&1&0&l_1\\-s_3&0&c_3&l_2\\0&0&0&1\end{bmatrix},
  \label{eq:qk-3d-T}
\end{equation}
其中 $l_1=\mp l_{abad}$（右腿取负 / 左腿取正）、大腿偏置项 $-l_2$ 表小腿系平移；足端在 $\{3\}$ 系坐标 $\boldsymbol{p}_3=[0,0,l_3]^\top$，$l_3=-l_{knee}$\cite{unitree2023quadruped}。

\begin{note}[左右腿 abad 偏置符号约定]\label{note:qk-ocr-leg-sign}
$l_1=\mp l_{abad}$ 的左右腿符号必须按腿别取定：右侧腿（右前 / 右后）$l_1=-l_{abad}$、左侧腿（左前 / 左后）$l_1=+l_{abad}$\cite{unitree2023quadruped}。该符号决定 abad 偏置朝机身内侧还是外侧，取错则整条腿镜像。装配 FK 时务必按腿别取号。
\end{note}

\paragraph{单腿正运动学闭式（逐级连乘，不跳步）。}
足端在 $\{3\}$ 系的齐次坐标是常量 $[\boldsymbol{p}_3;1]=[0,0,l_3,1]^\top$。把它依次左乘 $\mathbf{T}_{23}\to\mathbf{T}_{12}\to\mathbf{T}_{01}$，每一步把足端坐标“抬”到上一级坐标系，最终落到机身（腿根）系 $\{0\}$，即 $[\boldsymbol{p}_0;1]=\mathbf{T}_{01}\mathbf{T}_{12}\mathbf{T}_{23}[\boldsymbol{p}_3;1]$\cite{unitree2023quadruped}。逐级展开如下。

\emph{第一级}（$\{3\}\!\to\!\{2\}$）。$\mathbf{T}_{23}$ 第二、三列的偏置 $l_1,l_2$ 直接进入平移：
\begin{equation}
  \mathbf{T}_{23}\begin{bmatrix}0\\0\\l_3\\1\end{bmatrix}
  =\begin{bmatrix}c_3\!\cdot\!0+s_3 l_3\\[1pt] 0+l_1\\[1pt] -s_3\!\cdot\!0+c_3 l_3+l_2\\[1pt] 1\end{bmatrix}
  =\begin{bmatrix}l_3 s_3\\ l_1\\ l_3 c_3+l_2\\ 1\end{bmatrix}.
  \label{eq:qk-fk-step1}
\end{equation}
（几何：小腿绕 $y$ 转 $\theta_3$，把 $l_3$ 投到 $x,z$；再叠上大腿系偏置 $l_1$（$y$）、$l_2$（$z$）。）

\emph{第二级}（$\{2\}\!\to\!\{1\}$，绕 $y$ 转 $\theta_2$）：
\begin{equation}
  \mathbf{T}_{12}\begin{bmatrix}l_3 s_3\\ l_1\\ l_3 c_3+l_2\\ 1\end{bmatrix}
  =\begin{bmatrix}c_2 l_3 s_3+s_2(l_3 c_3+l_2)\\ l_1\\ -s_2 l_3 s_3+c_2(l_3 c_3+l_2)\\ 1\end{bmatrix}
  =\begin{bmatrix}l_3\sin(\theta_2{+}\theta_3)+l_2 s_2\\ l_1\\ l_3\cos(\theta_2{+}\theta_3)+l_2 c_2\\ 1\end{bmatrix},
  \label{eq:qk-fk-step2}
\end{equation}
末步用了和角公式 $s_3 c_2+c_3 s_2=\sin(\theta_2{+}\theta_3)$、$c_3 c_2-s_3 s_2=\cos(\theta_2{+}\theta_3)$。

\emph{第三级}（$\{1\}\!\to\!\{0\}$，绕 $x$ 转 $\theta_1$）。记 $a:=l_3\sin(\theta_2{+}\theta_3)+l_2 s_2$、$d:=l_3\cos(\theta_2{+}\theta_3)+l_2 c_2$，$\mathbf{T}_{01}$ 只混合 $y,z$ 分量：
\begin{equation}
  \mathbf{T}_{01}\begin{bmatrix}a\\ l_1\\ d\\ 1\end{bmatrix}
  =\begin{bmatrix}a\\ c_1 l_1-s_1 d\\ s_1 l_1+c_1 d\\ 1\end{bmatrix}.
  \label{eq:qk-fk-step3}
\end{equation}
把 $a,d$ 代回 \cref{eq:qk-fk-step3}，即得足端在机身（腿根）系的正运动学闭式\cite{unitree2023quadruped}：
\begin{equation}
  \begin{bmatrix}x_P\\y_P\\z_P\end{bmatrix}=
  \begin{bmatrix}
    l_3\sin(\theta_2+\theta_3)+l_2\sin\theta_2\\
    -l_3\sin\theta_1\cos(\theta_2+\theta_3)+l_1\cos\theta_1-l_2\cos\theta_2\sin\theta_1\\
    l_3\cos\theta_1\cos(\theta_2+\theta_3)+l_1\sin\theta_1+l_2\cos\theta_1\cos\theta_2
  \end{bmatrix}.
  \label{eq:qk-3d-fk}
\end{equation}
（核对：$y_P=c_1 l_1-s_1 d=l_1 c_1-s_1[l_3\cos(\theta_2{+}\theta_3)+l_2 c_2]$ 展开即第二行；$z_P=s_1 l_1+c_1 d$ 即第三行。）这三式是单腿运动学的“正算”核心：给定 $(\theta_1,\theta_2,\theta_3)$ 直接代入即得足端位置，唯一且无需迭代——正是 \cref{ins:qk-fk-easy-ik-hard} 中“FK 易”的体现。

\begin{example}[单腿正运动学数值例]\label{exam:qk-fk-num}
取 \cref{exam:qk-dims} 参数 $l_{abad}=0.088,\ l_{hip}=l_{knee}=0.27\,\mathrm{m}$，关节角 $\theta_1=0$、$\theta_2=-\tfrac{\pi}{4}$、$\theta_3=\tfrac{\pi}{2}$。务必按 \cref{eq:qk-3d-T} 与 \cref{note:qk-ocr-leg-sign} 的符号约定取连杆参数：$l_1=+0.088$（右腿本应 $-l_{abad}=-0.088$，此处取正以简化演示）、$l_2=-l_{hip}=-0.27$、$l_3=-l_{knee}=-0.27$（$l_2,l_3$ 取负是 \cref{eq:qk-3d-T} 坐标系约定下的硬性要求，漏负号会使足端 $z$ 反号）。则 $\theta_2+\theta_3=\tfrac{\pi}{4}$，$\sin(\theta_2+\theta_3)=\cos(\theta_2+\theta_3)=\tfrac{\sqrt2}{2}\approx0.707$，$\sin\theta_2=-0.707,\cos\theta_2=0.707$。代入 \cref{eq:qk-3d-fk}：
\[
  x_P=(-0.27)(0.707)+(-0.27)(-0.707)=0\,\mathrm{m};
\]
\[
  y_P=-(-0.27)\cdot0\cdot0.707+0.088\cdot1-(-0.27)\cdot0.707\cdot0=0.088\,\mathrm{m};
\]
\[
  z_P=(-0.27)\cdot1\cdot0.707+0.088\cdot0+(-0.27)\cdot1\cdot0.707=-0.382\,\mathrm{m}.
\]
即足端位于 $(0,\,0.088,\,-0.382)\,\mathrm{m}$：$z$ 轴朝上，故足端落在腿根\emph{正下方}（$z<0$）、沿 $y$ 偏外一点——大腿后摆、小腿前折的自然站姿。此结果将在 \cref{exam:qk-ik-num} 被逆运动学回代\emph{精确}复原 $\theta_1=0$，形成闭环自洽。
\end{example}

% ════════════════════════════════════════════════════════════════
\section{单腿逆运动学：解耦三步}
\label{sec:qk-ik}

逆运动学要由足端位置 $(x_P,y_P,z_P)$ 反解 $(\theta_1,\theta_2,\theta_3)$。一套\emph{解耦}的解析三步法是：先用投影几何单独解出 $\theta_1$，再用余弦定理解 $\theta_3$（不依赖 $\theta_2$），最后整理解出 $\theta_2$\cite{unitree2023quadruped}。这是单腿 IK 可\emph{闭式求解}（无需数值迭代）的关键。

\paragraph{第一步：$\theta_1$（yz 平面投影 $+$ atan2）。}
abad 关节绕 $x$ 轴，只影响 $yz$ 平面内的投影。把腿投影到 $yz$ 平面：足端在 $\{0\}$ 系的投影是 $(y_P,z_P)$，在随 abad 转过 $\theta_1$ 的 $\{1\}$ 系则是常量 $(l_1,-L)$，其中 $L=\sqrt{y_P^2+z_P^2-l_1^2}$ 为大、小腿在该平面的投影长（勾股：$y_P^2+z_P^2=l_1^2+L^2$）。二者由平面旋转相联，展开得\cite{unitree2023quadruped}
\begin{equation}
  \begin{bmatrix}y_P\\z_P\end{bmatrix}
  =\begin{bmatrix}c_1&-s_1\\ s_1&c_1\end{bmatrix}\begin{bmatrix}l_1\\ -L\end{bmatrix}
  \;\Longrightarrow\;
  \begin{cases} y_P=l_1 c_1+L s_1,\\[1pt] z_P=l_1 s_1-L c_1.\end{cases}
  \label{eq:qk-ik-th1-expand}
\end{equation}
两式同除 $c_1$ 得 $y_P/c_1=l_1+L\tan\theta_1$、$z_P/c_1=l_1\tan\theta_1-L$；再两式相除消去 $c_1$：
\[
  \frac{y_P}{z_P}=\frac{l_1+L\tan\theta_1}{l_1\tan\theta_1-L}
  \;\Longrightarrow\;
  y_P l_1\tan\theta_1-y_P L=z_P l_1+z_P L\tan\theta_1 .
\]
把含 $\tan\theta_1$ 的项归并、常数项移到右端，$\tan\theta_1\,(y_P l_1-z_P L)=z_P l_1+y_P L$，于是
\begin{equation}
  L=\sqrt{y_P^2+z_P^2-l_1^2},
  \qquad
  \theta_1=\operatorname{atan2}\!\big(z_P l_1+y_P L,\ y_P l_1-z_P L\big).
  \label{eq:qk-ik-theta1}
\end{equation}
用 atan2 而非 $\arctan$ 的单一比值，是为按分子、分母的符号唯一定出象限（\cref{note:qk-ocr-atan2}）。

\paragraph{第二步：$\theta_3$（余弦定理，不依赖 $\theta_2$）。}
设髋点 $O_3$、足端 $P$。大腿小腿在其平面内构成一个三角形 $\triangle O_3 A P$，由余弦定理直接解膝角\cite{unitree2023quadruped}：
\begin{equation}
  |\overrightarrow{AP}|=\sqrt{x_P^2+y_P^2+z_P^2-l_{abad}^2},
  \qquad
  \theta_3=-\pi+\arccos\!\Big(\frac{|\overrightarrow{O_3A}|^2+|\overrightarrow{O_3P}|^2-|\overrightarrow{AP}|^2}{2\,|\overrightarrow{O_3A}|\,|\overrightarrow{O_3P}|}\Big),
  \label{eq:qk-ik-theta3}
\end{equation}
其中 $|\overrightarrow{O_3A}|=l_{hip}$、$|\overrightarrow{O_3P}|=l_{knee}$。先算 $\theta_3$ 是因为它\emph{只}由三段长度的几何关系决定，与 $\theta_2$ 无关——这正是“解耦”的精髓。

\paragraph{第三步：$\theta_2$（整理 tan 分式 $+$ atan2）。}
已知 $\theta_1,\theta_3$ 后，\cref{eq:qk-3d-fk} 中仅 $\theta_2$ 未知。由 $x_P$ 式与 $(y_P\sin\theta_1-z_P\cos\theta_1)$ 式相除得 $\tan\theta_2$\cite{unitree2023quadruped}：
\begin{equation}
  \theta_2=\operatorname{atan2}\big(a_1 m_1+a_2 m_2,\ a_2 m_1-a_1 m_2\big),
  \quad
  \begin{cases}
    a_1=y_P\sin\theta_1-z_P\cos\theta_1,\\
    a_2=x_P,\\
    m_1=l_3\sin\theta_3,\\
    m_2=l_3\cos\theta_3+l_2.
  \end{cases}
  \label{eq:qk-ik-theta2}
\end{equation}
在样机限位范围内可证 $a_1,a_2$ 不同时为零，故 atan2 总有定义\cite{unitree2023quadruped}。

\begin{example}[单腿逆运动学数值例（回代验证）]\label{exam:qk-ik-num}
取 \cref{exam:qk-fk-num} 算出的足端 $(x_P,y_P,z_P)=(0,\,0.088,\,-0.382)$，投影偏置 $l_1=+0.088$（与上例同号）。
第一步：$L=\sqrt{0.088^2+(-0.382)^2-0.088^2}=0.382$；代入 \cref{eq:qk-ik-theta1}，
\[
  \theta_1=\operatorname{atan2}\!\big(z_P l_1+y_P L,\ y_P l_1-z_P L\big)
  =\operatorname{atan2}\!\big(\underbrace{(-0.382)(0.088)+(0.088)(0.382)}_{=0},\ \underbrace{(0.088)(0.088)-(-0.382)(0.382)}_{=0.1536>0}\big)=0.
\]
分子 $=-0.0336+0.0336=0$、分母 $=0.00774+0.1459>0$，故 $\theta_1=\operatorname{atan2}(0,+)=0$——\emph{精确}复原原 $\theta_1=0$（这正是上例采用 $l_2,l_3<0$ 符号约定的回报：$z_P<0$ 使分子两项恰好抵消）。
第二步：$|\overrightarrow{AP}|=\sqrt{0+0.0077+0.1459-0.0077}=\sqrt{0.1459}=0.382$；余弦项 $=\frac{0.27^2+0.27^2-0.382^2}{2\cdot0.27\cdot0.27}=\frac{0.0729+0.0729-0.1459}{0.1458}=\frac{-0.0001}{0.1458}\approx0$，$\arccos(0)=\tfrac{\pi}{2}$，故 $\theta_3=-\pi+\tfrac{\pi}{2}=-\tfrac{\pi}{2}$。
（注：此处与原 $\theta_3=+\tfrac{\pi}{2}$ 差一符号，对应膝关节“前折 / 后折”两个 IK 解，见下 \cref{pit:qk-ik-multi}。）回代验证 FK 自洽：取 $\theta_3=\tfrac{\pi}{2}$ 这一支即复原 \cref{exam:qk-fk-num}。此例展示 IK 闭式三步可逆地复原关节角，并自然暴露多解结构。
\end{example}

\begin{pitfall}[陷阱：IK 多解、无解与限位取舍]\label{pit:qk-ik-multi}
单腿 IK 天然存在\emph{多解}：abad 角 $\theta_1$ 理论上有两解（投影到 $yz$ 平面的两个对称构型），膝角 $\theta_3$ 也有“前折 / 后折”两支（\cref{exam:qk-ik-num} 的 $\pm\tfrac{\pi}{2}$）。工程上靠\textbf{关节限位}消歧：以 A1 为例，$\theta_1$ 的另一解超出物理限位故\emph{舍弃}，膝关节按机械结构只允许一个折向\cite{unitree2023quadruped}。此外若目标点超出可达工作空间（\cref{ins:qk-fk-easy-ik-hard}）则\emph{无解}，代码中表现为 $\arccos$ 的参数超出 $[-1,1]$——必须先做边界裁剪或可达性判断，否则会得到 NaN。这正是“IK 难”的工程体现。
\end{pitfall}

% ════════════════════════════════════════════════════════════════
\section{一阶微分运动学与静力学对偶：足端雅可比}
\label{sec:qk-jacobian}

\begin{insight}[K5：一阶微分把非线性变线性（全章方法论核心）]\label{ins:qk-diff-linearize}
这是贯穿全章、也最重要的方法论洞见。\cref{eq:qk-3d-fk} 表明足端位置 $\boldsymbol{p}$ 与关节角 $\boldsymbol{q}$ 是\emph{非线性}关系，\textbf{不能}写成矩阵乘 $\boldsymbol{p}=\mathbf{A}\boldsymbol{q}$——这正是 IK 难、需逐腿解耦的根源。但关键在于：\textbf{在\emph{速度}层面，二者却是\emph{线性}的}\cite{unitree2023quadruped}：
\[
  \dot{\boldsymbol{p}}=\frac{\partial\boldsymbol{f}}{\partial\boldsymbol{q}}\dot{\boldsymbol{q}}=\mathbf{J}(\boldsymbol{q})\,\dot{\boldsymbol{q}}.
\]
对固定的构型 $\boldsymbol{q}$，雅可比 $\mathbf{J}(\boldsymbol{q})$ 是一个\emph{常矩阵}，$\dot{\boldsymbol{p}}$ 与 $\dot{\boldsymbol{q}}$ 成线性关系。“线性关系是非常有用的属性”——因为线性映射可\emph{求逆}：$\dot{\boldsymbol{q}}=\mathbf{J}^{-1}\dot{\boldsymbol{p}}$（$\mathbf{J}$ 方阵可逆时），\emph{无需}解非线性 IK 即可在速度层面正逆切换。整个微分运动学、雅可比、零空间、WBC 的优化求解，本质都建立在“一阶微分把非线性变线性”这一招上。把它吃透，后面一切都顺。
\end{insight}

\begin{note}[关节角记号 $\theta_i$ 与 $q_i$]\label{note:qk-theta-q}
单腿闭式一节用 $\theta_i$ 表关节角、$\mathbf{J}$ 表单腿 $3\times3$ 足端雅可比；\cref{sec:qk-plucker} 起的整机框架用 $q_i$ 表关节变量、$\mathbf{J}_k$ 表整机雅可比。$\theta_i$ 即 $\boldsymbol{q}$ 的对应腿关节分量 $q_i$，所指相同，记号差异仅为贴合各自的维度与上下文。
\end{note}

\subsection{足端雅可比（FK 直接求导）}
\label{subsec:qk-jacobian-def}

把 \cref{eq:qk-3d-fk} 视为向量函数 $\boldsymbol{p}=\boldsymbol{f}(\boldsymbol{q})$，对关节角求偏导得足端雅可比\cite{unitree2023quadruped}（求导机制见 \cref{ch:matderiv}）：
\begin{equation}
  \dot{\boldsymbol{p}}=\begin{bmatrix}\dot x_P\\\dot y_P\\\dot z_P\end{bmatrix}
  =\frac{\partial\boldsymbol{f}}{\partial\boldsymbol{q}}\dot{\boldsymbol{q}}=\mathbf{J}(\boldsymbol{q})\,\dot{\boldsymbol{q}},
  \qquad
  \mathbf{J}=\frac{\partial\boldsymbol{f}}{\partial\boldsymbol{q}}\in\mathbb{R}^{3\times3}.
  \label{eq:qk-jacobian-def}
\end{equation}
逐元对 $\theta_1,\theta_2,\theta_3$ 求偏导，装配 $3\times3$ 雅可比。记 $c_{23}{=}\cos(\theta_2{+}\theta_3)$、$s_{23}{=}\sin(\theta_2{+}\theta_3)$，并用 $\partial c_{23}/\partial\theta_2=\partial c_{23}/\partial\theta_3=-s_{23}$、$\partial s_{23}/\partial\theta_j=c_{23}$、$\partial c_2/\partial\theta_2=-s_2$。

\emph{第一行}（$x_P=l_3 s_{23}+l_2 s_2$，不含 $\theta_1$，故 $J_{11}=0$）：
\begin{equation}
  \frac{\partial x_P}{\partial\theta_1}=0,\quad
  \frac{\partial x_P}{\partial\theta_2}=l_3 c_{23}+l_2 c_2,\quad
  \frac{\partial x_P}{\partial\theta_3}=l_3 c_{23}.
  \label{eq:qk-jacobian-row1}
\end{equation}

\emph{第二行}（$y_P=-l_3 s_1 c_{23}+l_1 c_1-l_2 c_2 s_1$）：
\begin{equation}
  \begin{aligned}
  \frac{\partial y_P}{\partial\theta_1}&=-l_3 c_1 c_{23}-l_1 s_1-l_2 c_1 c_2,\\
  \frac{\partial y_P}{\partial\theta_2}&=l_3 s_1 s_{23}+l_2 s_1 s_2=s_1\,(l_3 s_{23}+l_2 s_2),\\
  \frac{\partial y_P}{\partial\theta_3}&=l_3 s_1 s_{23}.
  \end{aligned}
  \label{eq:qk-jacobian-row2}
\end{equation}

\emph{第三行}（$z_P=l_3 c_1 c_{23}+l_1 s_1+l_2 c_1 c_2$）：
\begin{equation}
  \begin{aligned}
  \frac{\partial z_P}{\partial\theta_1}&=-l_3 s_1 c_{23}+l_1 c_1-l_2 s_1 c_2,\\
  \frac{\partial z_P}{\partial\theta_2}&=-l_3 c_1 s_{23}-l_2 c_1 s_2=-c_1\,(l_3 s_{23}+l_2 s_2),\\
  \frac{\partial z_P}{\partial\theta_3}&=-l_3 c_1 s_{23}.
  \end{aligned}
  \label{eq:qk-jacobian-row3}
\end{equation}

合成足端雅可比（$c_i{=}\cos\theta_i,\ s_i{=}\sin\theta_i$）\cite{unitree2023quadruped}：
\begin{equation}
  \mathbf{J}=\begin{bmatrix}
  0 & l_3 c_{23}+l_2 c_2 & l_3 c_{23}\\[2pt]
  -l_3 c_1 c_{23}-l_1 s_1-l_2 c_1 c_2 & s_1(l_3 s_{23}+l_2 s_2) & l_3 s_1 s_{23}\\[2pt]
  -l_3 s_1 c_{23}+l_1 c_1-l_2 s_1 c_2 & -c_1(l_3 s_{23}+l_2 s_2) & -l_3 c_1 s_{23}
  \end{bmatrix}.
  \label{eq:qk-jacobian-full}
\end{equation}

\paragraph{速度层正逆。}
在 A1/Go1 的限位范围内 $\mathbf{J}$ 可逆，故由 \cref{ins:qk-diff-linearize} 可在速度层正逆切换\cite{unitree2023quadruped}：
\begin{equation}
  \dot{\boldsymbol{p}}=\mathbf{J}\dot{\boldsymbol{q}}\quad(\text{正}),
  \qquad
  \dot{\boldsymbol{q}}=\mathbf{J}^{-1}\dot{\boldsymbol{p}}\quad(\text{逆}).
  \label{eq:qk-jacobian-invert}
\end{equation}

\subsection{静力学对偶：\texorpdfstring{$\boldsymbol{\tau}=\mathbf{J}^\top\mathbf{F}$}{tau=J^TF}}
\label{subsec:qk-statics}

\begin{insight}[K6：雅可比的运动学 / 静力学对偶（功率守恒）]\label{ins:qk-duality}
一个深刻而实用的事实是：\textbf{同一个雅可比 $\mathbf{J}$，既贯通速度映射 $\dot{\boldsymbol{p}}=\mathbf{J}\dot{\boldsymbol{q}}$，又贯通力映射 $\boldsymbol{\tau}=\mathbf{J}^\top\mathbf{F}$}。其物理根源是\emph{功率守恒}：当整条腿处于（准）静止平衡时，关节侧输入的功率必等于足端侧输出的功率（流入功率 $=$ 流出功率，无功率储存于无质量连杆）\cite{unitree2023quadruped}。这就是“运动学 / 静力学对偶”——速度与力是一对偶量，共用一个 $\mathbf{J}$（一个用 $\mathbf{J}$、一个用 $\mathbf{J}^\top$）。在 \cref{ch:quad_mpc} 里，MPC 算出足端力 $\mathbf{f}_i$ 后，正是用这条对偶把它转成电机力矩；在 \cref{ch:quad_wbc} 里，它是把任务空间力映射回关节空间的桥。
\end{insight}

\begin{theorem}[运动学 / 静力学对偶]\label{thm:qk-duality}
设足端雅可比 $\mathbf{J}$ 满足 $\dot{\boldsymbol{p}}=\mathbf{J}\dot{\boldsymbol{q}}$。若整条腿处于准静态平衡，则足端对外力 $\mathbf{F}$ 与关节力矩 $\boldsymbol{\tau}$ 之间满足 $\boldsymbol{\tau}=\mathbf{J}^\top\mathbf{F}$。
\end{theorem}
\begin{proof}
准静态下功率守恒要求关节侧功率等于足端侧功率，对任意（虚）速度成立\cite{unitree2023quadruped}：
\begin{equation}
  \tau_1\dot\theta_1+\tau_2\dot\theta_2+\tau_3\dot\theta_3
  =F_x\dot x_P+F_y\dot y_P+F_z\dot z_P
  \quad\Longleftrightarrow\quad
  \boldsymbol{\tau}^\top\dot{\boldsymbol{q}}=\mathbf{F}^\top\dot{\boldsymbol{p}}.
  \label{eq:qk-power-balance}
\end{equation}
代入 $\dot{\boldsymbol{p}}=\mathbf{J}\dot{\boldsymbol{q}}$：
\[
  \boldsymbol{\tau}^\top\dot{\boldsymbol{q}}=\mathbf{F}^\top\mathbf{J}\dot{\boldsymbol{q}}
  =(\mathbf{J}^\top\mathbf{F})^\top\dot{\boldsymbol{q}}.
\]
由于该式对\emph{任意} $\dot{\boldsymbol{q}}$ 成立，两边系数相等，得 $\boldsymbol{\tau}=\mathbf{J}^\top\mathbf{F}$。\qed
\end{proof}

\paragraph{正逆与符号约定。}
正逆映射为 $\boldsymbol{\tau}=\mathbf{J}^\top\mathbf{F}$、$\mathbf{F}=\mathbf{J}^{-\top}\boldsymbol{\tau}$。须注意力的方向约定：此处 $\mathbf{F}$ 取为足端\emph{对地}力；若以地对足的反作用力 $\mathbf{F}'$ 表达，则 $\boldsymbol{\tau}=-\mathbf{J}^\top\mathbf{F}'$\cite{unitree2023quadruped}。这一符号在 \cref{ch:quad_mpc} 把 GRF 转关节力矩时必须对齐。

\begin{practice}[摆动腿笛卡尔阻抗（工程，凝练不全推）]\label{prac:qk-swing}
支撑腿用 \cref{thm:qk-duality} 把规划的足端力转关节力矩；\emph{摆动腿}则用笛卡尔空间 PD（阻抗）跟踪期望足端轨迹，再经同一对偶转力矩\cite{unitree2023quadruped}：
\[
  \mathbf{f}_d=\mathbf{K}_p(\boldsymbol{p}_{0d}-\boldsymbol{p}_{0f})+\mathbf{K}_d(\dot{\boldsymbol{p}}_{0d}-\dot{\boldsymbol{p}}_{0f}),
  \qquad \boldsymbol{\tau}=\mathbf{J}^\top\mathbf{f}_d,
\]
其中 $\mathbf{K}_p,\mathbf{K}_d$ 为正定对角增益。这属于工程实现细节，本书凝练于此、不展开整定；轨迹生成见 \cref{ch:quad_gait}，与 MPC 的力分配协同见 \cref{ch:quad_wbc}。
\end{practice}

% ════════════════════════════════════════════════════════════════
\section{Plücker 空间向量与递推正运动学}
\label{sec:qk-plucker}

单腿的齐次变换法直观、可手算，但难以系统地推广到“$13$ 刚体 $+$ 浮动基座 $+$ 速度 / 加速度 / 力同时递推”。本节引入 Featherstone 的\textbf{Plücker 六维空间向量}代数\cite{yu2021quadruped,featherstone2008rigid}，它把运动与力统一成六维向量，配运动学树给出可任意扩展的递推算法。这是现代多刚体动力学的统一语言，也是 \cref{ch:quad_wbc} 全身控制的底层框架。

\subsection{六维运动 / 力向量}
\label{subsec:qk-spatial-vec}

\begin{insight}[为何要把平动与转动“合一”：传统两套算法的割裂]\label{ins:qk-why-spatial}
引入六维空间向量的根本动机是：\textbf{传统刚体动力学把运动割裂为“$3$ 维平动 $+$ $3$ 维转动”，企图用两套算法（牛顿第二定律管平动、欧拉方程管转动）分别处理；但平动与转动\emph{无法简单分离}}\cite{yu2021quadruped}——刚体上的点会\emph{因转动而获得线速度}（\cref{eq:qk-angvel-1} 的 $\boldsymbol{\omega}\times\boldsymbol{r}$ 项），一个力施加到刚体上\emph{既改变平动、又改变转动}。这种耦合渗透在动力学算法的每一处，使两套算法既不简洁、又难推导。Plücker 六维向量把“转动 $\boldsymbol{\omega}$”与“平动 $\boldsymbol{v}_O$”塞进\emph{同一个}六维对象，让坐标变换、叉乘、运动方程都成为\emph{一条}规则——这正是 \cref{note:qk-modeling} 中“第三代建模范式”相对前两代的根本优势。代价是定义略繁，但一旦接受，整机递推（\cref{alg:qk-recursive-fk}）与全身雅可比（\cref{sec:qk-wbjac}）都变得统一可扩展。
\end{insight}

\paragraph{空间速度（运动向量）与 Plücker 基。}
回到 \cref{eq:qk-angvel-1}：刚体的速度场被“纯转动 $\boldsymbol{\omega}$（绕过参考点 $O$ 的某直线）$+$ 过 $O$ 的纯平动 $\boldsymbol{v}_O$”完整刻画——给定刚体上任一点位置即可由这两者算出其速度，故二者已是\emph{完整}描述\cite{yu2021quadruped}。在 $O$ 处建笛卡尔单位正交基 $\{\boldsymbol{i},\boldsymbol{j},\boldsymbol{k}\}$，把 $\boldsymbol{\omega},\boldsymbol{v}_O$ 各自展开：
\begin{equation}
  \boldsymbol{\omega}=[\boldsymbol{i}\ \boldsymbol{j}\ \boldsymbol{k}]\,[\omega_x\ \omega_y\ \omega_z]^\top,
  \qquad
  \boldsymbol{v}_O=[\boldsymbol{i}\ \boldsymbol{j}\ \boldsymbol{k}]\,[v_{Ox}\ v_{Oy}\ v_{Oz}]^\top.
  \label{eq:qk-plucker-basis-3d}
\end{equation}
为把二者并入\emph{一个}向量，定义一组\emph{六维}正交基 $\mathcal{D}_O=\{\boldsymbol{d}_{Ox},\boldsymbol{d}_{Oy},\boldsymbol{d}_{Oz},\boldsymbol{d}_x,\boldsymbol{d}_y,\boldsymbol{d}_z\}\subseteq\mathrm{M}^6$，其中 $\boldsymbol{d}_{O\bullet}$ 是绕 $\bullet$ 轴的\emph{单位转动}、$\boldsymbol{d}_\bullet$ 是沿 $\bullet$ 轴的\emph{单位平动}。于是刚体空间速度展开为这六个基的线性组合，参数化即得六维 Plücker 运动向量\cite{yu2021quadruped}：
\begin{equation}
  \boldsymbol{v}=\omega_x\boldsymbol{d}_{Ox}+\omega_y\boldsymbol{d}_{Oy}+\omega_z\boldsymbol{d}_{Oz}+v_{Ox}\boldsymbol{d}_x+v_{Oy}\boldsymbol{d}_y+v_{Oz}\boldsymbol{d}_z
  =[\,\omega_x\ \omega_y\ \omega_z\ \ v_{Ox}\ v_{Oy}\ v_{Oz}\,]^\top=[\,\boldsymbol{\omega}^\top\ \boldsymbol{v}_O^\top\,]^\top.
  \label{eq:qk-spatial-vel}
\end{equation}
$\mathrm{M}^6$ 称\textbf{运动向量空间}（motion space），其上可定义空间速度、空间加速度等；“一个 Plücker 坐标系依赖于一个笛卡尔坐标系，每个向量须用 $6$ 个参数描述”\cite{yu2021quadruped}。本书以前缀“空间”表“六维向量”或“$6\times6$ 张量”。

\paragraph{空间力（力向量）与对偶基。}
对偶地，刚体所受力对\emph{过点 $P$ 的某直线上的力 $\boldsymbol{f}$ $+$ 该点力矩 $\boldsymbol{n}_P$}分解；换参考点 $O$ 时力 $\boldsymbol{f}$ 不变，力矩按\emph{力矩随参考点平移}的标准关系变换\cite{yu2021quadruped}：
\begin{equation}
  \boldsymbol{n}_P=\boldsymbol{n}_O+\boldsymbol{f}\times\overrightarrow{OP}.
  \label{eq:qk-wrench-shift}
\end{equation}
在 $O$ 处定义六维正交基 $\mathcal{E}_O=\{\boldsymbol{e}_x,\boldsymbol{e}_y,\boldsymbol{e}_z,\boldsymbol{e}_{Ox},\boldsymbol{e}_{Oy},\boldsymbol{e}_{Oz}\}\subseteq\mathrm{F}^6$（$\boldsymbol{e}_\bullet$ 绕 $\bullet$ 轴单位力矩、$\boldsymbol{e}_{O\bullet}$ 沿 $\bullet$ 轴单位力），展开并参数化得六维 Plücker 力向量\cite{yu2021quadruped}：
\begin{equation}
  \hat{\boldsymbol{f}}=n_{Ox}\boldsymbol{e}_x+n_{Oy}\boldsymbol{e}_y+n_{Oz}\boldsymbol{e}_z+f_x\boldsymbol{e}_{Ox}+f_y\boldsymbol{e}_{Oy}+f_z\boldsymbol{e}_{Oz}
  =[\,n_{Ox}\ n_{Oy}\ n_{Oz}\ \ f_x\ f_y\ f_z\,]^\top=[\,\boldsymbol{n}_O^\top\ \boldsymbol{f}^\top\,]^\top.
  \label{eq:qk-spatial-force}
\end{equation}
$\mathrm{F}^6$ 称\textbf{动力向量空间}（force space），其上可定义空间力、空间动量。当六维空间向量与传统三维向量有符号冲突时，在六维量上加“尖帽”$\hat{(\cdot)}$。\textbf{关键}：尽管三维下力、力矩、速度、加速度同处一个坐标系、同套运算，六维下\emph{运动向量空间与动力向量空间的运算规则却不同}（叉乘分 $\times$ 与 $\times^{*}$ 两套、变换分 $\mathbf{X}$ 与 $\mathbf{X}^{*}$ 两套，见下）——这是空间向量代数最需留心处。

\begin{note}[Plücker 分块顺序：本节角量在前，与全书 se(3) 位姿坐标的换算]\label{note:qk-plucker-order}
本节 Plücker 递推内部，空间向量采用 \textbf{角量在前}：运动向量取 $[\boldsymbol{\omega}^\top,\boldsymbol{v}_O^\top]^\top$、力向量取 $[\boldsymbol{n}_O^\top,\boldsymbol{f}^\top]^\top$（Featherstone 空间向量代数的标准分块约定）。于是六维向量上 $3$ 块恒为角量、下 $3$ 块恒为线量，下文 $\mathbf{X}$、$\boldsymbol{v}\times$、$\hat{\mathbf{I}}$ 等 $6\times6$ 矩阵的分块结构均按此排序给出。全书 $\se(3)$ \emph{位姿坐标} 则用\emph{平移在前} $\boldsymbol{\xi}=[\boldsymbol{\rho}^\top,\boldsymbol{\phi}^\top]^\top$（\cref{note:qk-twist-order}）；二者仅差上下 $3$ 维分块的交换，对应的 $\mathbf{X}$ 按分块转置交换即可换算。跨章对接 \cref{ch:quad_wbc} 的 CRBA/RNEA 时，$\mathbf{X}/\mathbf{X}^{*}/\mathbf{S}$ 与本节同套排序，可直接复用。
\end{note}

\subsection{Plücker 坐标变换（运动 / 力对偶）}
\label{subsec:qk-plucker-transform}

设两系间旋转 $\mathbf{E}={}^{\mathcal{B}}\mathbf{R}_{\mathcal{A}}$、平移 $\boldsymbol{r}=\overrightarrow{AB}$（在 $\mathcal{A}$ 系）。运动向量与力向量的坐标变换为\cite{yu2021quadruped}：
\begin{equation}
  {}^{\mathcal{B}}\boldsymbol{m}={}^{\mathcal{B}}\mathbf{X}_{\mathcal{A}}\,{}^{\mathcal{A}}\boldsymbol{m},
  \qquad
  {}^{\mathcal{B}}\hat{\boldsymbol{f}}={}^{\mathcal{B}}\mathbf{X}_{\mathcal{A}}^{*}\,{}^{\mathcal{A}}\hat{\boldsymbol{f}},
  \qquad
  {}^{\mathcal{B}}\mathbf{X}_{\mathcal{A}}^{*}={}^{\mathcal{B}}\mathbf{X}_{\mathcal{A}}^{-\top},
  \label{eq:qk-plucker-X}
\end{equation}
其中力变换 $\mathbf{X}^{*}=\mathbf{X}^{-\top}$ 由\emph{保功率不变} $\boldsymbol{m}^\top\hat{\boldsymbol{f}}$ 在两系相等导出（运动与力是对偶量）。两矩阵的显式分块为\cite{yu2021quadruped}：
\begin{equation}
  {}^{\mathcal{B}}\mathbf{X}_{\mathcal{A}}=\begin{bmatrix}\mathbf{E}&\mathbf{0}\\-\mathbf{E}\,\boldsymbol{r}^{\wedge}&\mathbf{E}\end{bmatrix},
  \qquad
  {}^{\mathcal{B}}\mathbf{X}_{\mathcal{A}}^{*}=\begin{bmatrix}\mathbf{E}&-\mathbf{E}\,\boldsymbol{r}^{\wedge}\\\mathbf{0}&\mathbf{E}\end{bmatrix}.
  \label{eq:qk-plucker-X-block}
\end{equation}
可分解为纯旋转 $\operatorname{diag}(\mathbf{E},\mathbf{E})$ 与纯平移 $\begin{bsmallmatrix}\mathbf{I}&\mathbf{0}\\-\boldsymbol{r}^{\wedge}&\mathbf{I}\end{bsmallmatrix}$ 的复合\cite{yu2021quadruped}。下面把这两个基本变换各自推一遍，再组合。

\paragraph{基本变换之一：纯旋转。}
设 $\mathcal{A},\mathcal{B}$ 共原点、仅差旋转 $\mathbf{E}={}^{\mathcal{B}}\mathbf{R}_{\mathcal{A}}$。六维运动向量 $\hat{\boldsymbol{m}}$ 的上、下三维块（角量 $\boldsymbol{m}$、线量 $\boldsymbol{m}_O$）\emph{各自}是普通三维向量，故各自左乘 $\mathbf{E}$ 换系，合写成块对角\cite{yu2021quadruped}：
\begin{equation}
  {}^{\mathcal{B}}\hat{\boldsymbol{m}}
  =\begin{bmatrix}\mathbf{E}\,{}^{\mathcal{A}}\boldsymbol{m}\\ \mathbf{E}\,{}^{\mathcal{A}}\boldsymbol{m}_O\end{bmatrix}
  =\begin{bmatrix}\mathbf{E}&\mathbf{0}\\\mathbf{0}&\mathbf{E}\end{bmatrix}{}^{\mathcal{A}}\hat{\boldsymbol{m}}
  \;\Rightarrow\;
  {}^{\mathcal{B}}\mathbf{X}_{\mathcal{A}}^{\text{rot}}=\begin{bmatrix}\mathbf{E}&\mathbf{0}\\\mathbf{0}&\mathbf{E}\end{bmatrix},
  \quad
  {}^{\mathcal{B}}\mathbf{X}_{\mathcal{A}}^{*\,\text{rot}}={}^{\mathcal{B}}\mathbf{X}_{\mathcal{A}}^{-\top}=\begin{bmatrix}\mathbf{E}&\mathbf{0}\\\mathbf{0}&\mathbf{E}\end{bmatrix}.
  \label{eq:qk-plucker-rot}
\end{equation}
（块对角正交阵的逆转置等于自身，故纯旋转下运动 / 力变换\emph{同形}。）

\paragraph{基本变换之二：纯平移。}
设 $\mathcal{A},\mathcal{B}$ 朝向相同、仅差平移 $\boldsymbol{r}=\overrightarrow{AB}$。角量 $\boldsymbol{m}$ 与参考点无关、不变；线量按 \cref{eq:qk-spatial-vel} 的速度场随参考点平移 $\boldsymbol{m}_B=\boldsymbol{m}_A-\overrightarrow{AB}\times\boldsymbol{m}=\boldsymbol{m}_A-\boldsymbol{r}^{\wedge}\boldsymbol{m}$\cite{yu2021quadruped}：
\begin{equation}
  {}^{\mathcal{B}}\hat{\boldsymbol{m}}
  =\begin{bmatrix}\boldsymbol{m}\\ \boldsymbol{m}_A-\boldsymbol{r}^{\wedge}\boldsymbol{m}\end{bmatrix}
  =\begin{bmatrix}\mathbf{I}&\mathbf{0}\\-\boldsymbol{r}^{\wedge}&\mathbf{I}\end{bmatrix}{}^{\mathcal{A}}\hat{\boldsymbol{m}}
  \;\Rightarrow\;
  {}^{\mathcal{B}}\mathbf{X}_{\mathcal{A}}^{\text{tr}}=\begin{bmatrix}\mathbf{I}&\mathbf{0}\\-\boldsymbol{r}^{\wedge}&\mathbf{I}\end{bmatrix},
  \quad
  {}^{\mathcal{B}}\mathbf{X}_{\mathcal{A}}^{*\,\text{tr}}={}^{\mathcal{B}}\mathbf{X}_{\mathcal{A}}^{-\top\,\text{tr}}=\begin{bmatrix}\mathbf{I}&-\boldsymbol{r}^{\wedge}\\\mathbf{0}&\mathbf{I}\end{bmatrix}.
  \label{eq:qk-plucker-tr}
\end{equation}
（对偶力变换的平移项移到\emph{右上}，恰对应 \cref{eq:qk-wrench-shift} 的力矩随点平移。）

\paragraph{组合变换与反向。}
先平移再旋转，复合即得 \cref{eq:qk-plucker-X-block}；连同反向变换 ${}^{\mathcal{A}}\mathbf{X}_{\mathcal{B}}={}^{\mathcal{B}}\mathbf{X}_{\mathcal{A}}^{-1}$ 一并列出（$\mathbf{E}^\top={}^{\mathcal{A}}\mathbf{R}_{\mathcal{B}}$）\cite{yu2021quadruped}：
\begin{equation}
  {}^{\mathcal{A}}\mathbf{X}_{\mathcal{B}}=\begin{bmatrix}\mathbf{E}^\top&\mathbf{0}\\\mathbf{E}^\top\boldsymbol{r}^{\wedge}&\mathbf{E}^\top\end{bmatrix},
  \qquad
  {}^{\mathcal{A}}\mathbf{X}_{\mathcal{B}}^{*}=\begin{bmatrix}\mathbf{E}^\top&\mathbf{E}^\top\boldsymbol{r}^{\wedge}\\\mathbf{0}&\mathbf{E}^\top\end{bmatrix}.
  \label{eq:qk-plucker-X-rev}
\end{equation}
平移项在反向变换里\emph{符号翻正}（$+\mathbf{E}^\top\boldsymbol{r}^{\wedge}$），与正向的 $-\mathbf{E}\boldsymbol{r}^{\wedge}$ 互逆——递推算法里频繁在父子系间来回搬运空间量（\cref{alg:qk-recursive-fk}），这组正反式是其底层运算。力变换 $\mathbf{X}^{*}=\mathbf{X}^{-\top}$ 的保功率来历见 \cref{deriv:qk-plucker-dual}。

\subsection{空间叉乘、空间加速度与空间惯量}
\label{subsec:qk-spatial-rest}

\paragraph{空间叉乘（运动 / 力两套）。}
三维下 \cref{eq:qk-angvel-1} 已表明角速度把刚接的常向量微分为 $\dot{\boldsymbol{r}}=\boldsymbol{\omega}\times\boldsymbol{r}$，即“$\boldsymbol{\omega}\times$”实现“向量$\to$其微分”的映射。这一性质推广到 Plücker 坐标系：设六维运动向量 $\boldsymbol{m}$、力向量 $\hat{\boldsymbol{f}}$ 为\emph{某坐标系下的常量}，该系在世界系下以空间速度 $\boldsymbol{v}$ 运动，则\cite{yu2021quadruped}
\begin{equation}
  {}^{\mathcal{O}}\dot{\boldsymbol{m}}=\boldsymbol{v}\times{}^{\mathcal{O}}\boldsymbol{m},
  \qquad
  {}^{\mathcal{O}}\dot{\hat{\boldsymbol{f}}}=\boldsymbol{v}\times^{*}{}^{\mathcal{O}}\hat{\boldsymbol{f}}.
  \label{eq:qk-spatial-cross-def}
\end{equation}
\emph{为何运动 / 力各需一套算子}？以基向量的微分定其矩阵元（逐基分析法）：若坐标系以单位角速度 $\boldsymbol{d}_{Ox}$ 绕 $Ox$ 旋转，经 $\delta t$ 后旋转轴 $Ox$ 不变（$\dot{\boldsymbol{d}}_{Ox}=\mathbf{0}$），但 $Oy$ 在 $y$-$z$ 平面转过 $\delta t$ 弧度，新系单位转动 $\boldsymbol{d}_{Oy}(t+\delta t)$ 在原系等于 $\boldsymbol{d}_{Oy}(t)+\delta t\,\boldsymbol{d}_{Oz}$，故 $\dot{\boldsymbol{d}}_{Oy}=\boldsymbol{d}_{Oz}$\cite{yu2021quadruped}：
\begin{equation}
  \dot{\boldsymbol{d}}_{Ox}=\lim_{\delta t\to0}\frac{\boldsymbol{d}_{Ox}(t+\delta t)-\boldsymbol{d}_{Ox}(t)}{\delta t}=\mathbf{0},
  \qquad
  \dot{\boldsymbol{d}}_{Oy}=\lim_{\delta t\to0}\frac{[\boldsymbol{d}_{Oy}(t)+\delta t\,\boldsymbol{d}_{Oz}]-\boldsymbol{d}_{Oy}(t)}{\delta t}=\boldsymbol{d}_{Oz}.
  \label{eq:qk-spatial-cross-micro}
\end{equation}
若坐标系改以单位平动 $\boldsymbol{d}_x$ 沿 $x$ 移动，则转动轴 $Oy$ 在该过程中于 $O$ 点产生一个沿 $z$ 的线速度，给出 $\dot{\boldsymbol{d}}_{Oy}=\boldsymbol{d}_z$（一个\emph{转动}基的微分落到\emph{平动}基上——这正是 $\boldsymbol{v}\times$ 左下出现 $\boldsymbol{v}_O^{\wedge}$ 的来由）。把六个单位运动逐一这样分析、对六运动六力共得 $72$ 个微分结果，排列即得两套算子\cite{yu2021quadruped}：
\begin{equation}
  \boldsymbol{v}\times=\begin{bmatrix}\boldsymbol{\omega}^{\wedge}&\mathbf{0}\\\boldsymbol{v}_O^{\wedge}&\boldsymbol{\omega}^{\wedge}\end{bmatrix},
  \qquad
  \boldsymbol{v}\times^{*}=\begin{bmatrix}\boldsymbol{\omega}^{\wedge}&\boldsymbol{v}_O^{\wedge}\\\mathbf{0}&\boldsymbol{\omega}^{\wedge}\end{bmatrix}=-(\boldsymbol{v}\times)^\top.
  \label{eq:qk-spatial-cross}
\end{equation}
$\boldsymbol{v}\times^{*}=-(\boldsymbol{v}\times)^\top$ 又一次体现运动 / 力的对偶（与 $\mathbf{X}^{*}=\mathbf{X}^{-\top}$ 同源）。

\paragraph{空间加速度。}
空间加速度定义为空间速度的时间导数 $\boldsymbol{a}=\dot{\boldsymbol{v}}=[\dot{\boldsymbol{\omega}}^\top\ \dot{\boldsymbol{v}}_O^\top]^\top$。看似显然，但 $\dot{\boldsymbol{v}}_O$ 的求法与传统加速度不同：$\boldsymbol{v}_O$ 是“\emph{当前}时刻与空间不动点 $O$ 重合的刚体点 $O'$ 的速度”，而 $O'$ 随刚体运动、$t$ 与 $t+\delta t$ 时刻与 $O$ 重合的是\emph{不同}刚体点。记 $\boldsymbol{r}$ 为 $O'$ 的位置：$t$ 时刻 $O'$ 与 $O$ 重合故 $\boldsymbol{v}_O(t)=\dot{\boldsymbol{r}}(t)$，但 $t+\delta t$ 时 $O'$ 已离开 $O$、走过 $\delta t\,\dot{\boldsymbol{r}}(t)$，故须\emph{扣除}这段位移随转动产生的速度\cite{yu2021quadruped}：
\begin{equation}
  \boldsymbol{v}_O(t+\delta t)=\dot{\boldsymbol{r}}(t+\delta t)-\boldsymbol{\omega}\times\delta t\,\dot{\boldsymbol{r}}(t)
  \;\Longrightarrow\;
  \dot{\boldsymbol{v}}_O=\lim_{\delta t\to0}\frac{\boldsymbol{v}_O(t+\delta t)-\boldsymbol{v}_O(t)}{\delta t}=\ddot{\boldsymbol{r}}-\boldsymbol{\omega}\times\dot{\boldsymbol{r}}.
  \label{eq:qk-spatial-acc-deriv}
\end{equation}
于是
\begin{equation}
  \boldsymbol{a}=\frac{\mathrm{d}\boldsymbol{v}}{\mathrm{d}t}
  =\begin{bmatrix}\dot{\boldsymbol{\omega}}\\\dot{\boldsymbol{v}}_O\end{bmatrix}
  =\begin{bmatrix}\dot{\boldsymbol{\omega}}\\\ddot{\boldsymbol{r}}-\boldsymbol{\omega}\times\dot{\boldsymbol{r}}\end{bmatrix}.
  \label{eq:qk-spatial-acc}
\end{equation}
\textbf{关键}：因 $O$ 是空间不动点、与之重合的刚体点逐刻不同，$\dot{\boldsymbol{v}}_O=\ddot{\boldsymbol{r}}-\boldsymbol{\omega}\times\dot{\boldsymbol{r}}$ \emph{不等于}“随体点加速度”$\ddot{\boldsymbol{r}}$——这一点是空间向量代数最易误解处\cite{yu2021quadruped}，也是 \cref{eq:qk-recursive} 加速度递推里 $\boldsymbol{v}_i\times\mathbf{S}_i\dot{\boldsymbol{q}}_i$ 项的物理根源。

\paragraph{空间动量与平行轴定理。}
设刚体质心 $C$ 当前与空间不动点重合、质量 $m$、质心惯性张量 $\bar{\mathbf{I}}_C$、$C$ 处空间速度 $\hat{\boldsymbol{v}}_C=[\boldsymbol{\omega}^\top\ \boldsymbol{v}_C^\top]^\top$。线动量 $\boldsymbol{h}=m\boldsymbol{v}_C$、固有角动量 $\boldsymbol{h}_C=\bar{\mathbf{I}}_C\boldsymbol{\omega}$；对任意点 $O$，动量矩 $=$ 固有角动量 $+$ 线动量的矩\cite{yu2021quadruped}：
\begin{equation}
  \boldsymbol{h}_O=\boldsymbol{h}_C+\overrightarrow{OC}\times\boldsymbol{h}
  \;\Longrightarrow\;
  \hat{\boldsymbol{h}}_O=\begin{bmatrix}\boldsymbol{h}_O\\\boldsymbol{h}\end{bmatrix}
  =\begin{bmatrix}\bar{\mathbf{I}}_C\boldsymbol{\omega}+\overrightarrow{OC}\times m\boldsymbol{v}_C\\ m\boldsymbol{v}_C\end{bmatrix}
  =\begin{bmatrix}\mathbf{I}_3&\boldsymbol{c}^{\wedge}\\\mathbf{0}&\mathbf{I}_3\end{bmatrix}\hat{\boldsymbol{h}}_C,
  \label{eq:qk-spatial-momentum}
\end{equation}
其中 $\boldsymbol{c}=\overrightarrow{OC}$。

\paragraph{空间惯量。}
仿线动量构建空间速度到空间动量的映射 $\hat{\boldsymbol{h}}=\hat{\mathbf{I}}\boldsymbol{v}$，$6\times6$ 矩阵 $\hat{\mathbf{I}}$ 即\emph{空间惯量}。质心处由 $\hat{\boldsymbol{h}}_C=[\bar{\mathbf{I}}_C\boldsymbol{\omega};\,m\boldsymbol{v}_C]=\operatorname{diag}(\bar{\mathbf{I}}_C,m\mathbf{I}_3)\hat{\boldsymbol{v}}_C$ 直接读出 $\hat{\mathbf{I}}_C$；一般点 $O$ 处把 \cref{eq:qk-spatial-momentum} 的平移与 $\hat{\boldsymbol{v}}_O={}\begin{bsmallmatrix}\mathbf{I}&\mathbf{0}\\\boldsymbol{c}^{\wedge\top}&\mathbf{I}\end{bsmallmatrix}\hat{\boldsymbol{v}}_C$ 夹乘（左乘平移、右乘其转置，平行轴定理的空间形式）\cite{yu2021quadruped}：
\begin{equation}
  \hat{\mathbf{I}}_C=\begin{bmatrix}\bar{\mathbf{I}}_C&\mathbf{0}\\\mathbf{0}&m\mathbf{I}_3\end{bmatrix},
  \qquad
  \hat{\mathbf{I}}_O=\begin{bmatrix}\mathbf{I}_3&\boldsymbol{c}^{\wedge}\\\mathbf{0}&\mathbf{I}_3\end{bmatrix}\hat{\mathbf{I}}_C\begin{bmatrix}\mathbf{I}_3&\mathbf{0}\\\boldsymbol{c}^{\wedge\top}&\mathbf{I}_3\end{bmatrix}
  =\begin{bmatrix}\bar{\mathbf{I}}_C+m\,\boldsymbol{c}^{\wedge}\boldsymbol{c}^{\wedge\top}&m\,\boldsymbol{c}^{\wedge}\\ m\,\boldsymbol{c}^{\wedge\top}&m\mathbf{I}_3\end{bmatrix}.
  \label{eq:qk-spatial-inertia}
\end{equation}
左上块 $\bar{\mathbf{I}}_C+m\boldsymbol{c}^{\wedge}\boldsymbol{c}^{\wedge\top}$ 即经典平行轴定理（惯量随参考点移而增 $m\boldsymbol{c}^{\wedge}\boldsymbol{c}^{\wedge\top}$）。

\paragraph{运动方程（牛顿--欧拉合一）。}
空间力等于空间动量的时间导数；用乘积法则 $\tfrac{\mathrm d}{\mathrm dt}(\hat{\mathbf{I}}\boldsymbol{v})=\hat{\mathbf{I}}\boldsymbol{a}+\dot{\hat{\mathbf{I}}}\boldsymbol{v}$，而 $\dot{\hat{\mathbf{I}}}$ 的贡献整理成 $(\boldsymbol{v}\times^{*}\hat{\mathbf{I}}-\hat{\mathbf{I}}\,\boldsymbol{v}\times)\boldsymbol{v}=\boldsymbol{v}\times^{*}\hat{\mathbf{I}}\boldsymbol{v}$（因 $\hat{\mathbf{I}}\,\boldsymbol{v}\times\boldsymbol{v}=\mathbf{0}$）\cite{yu2021quadruped}：
\begin{equation}
  \hat{\boldsymbol{f}}=\frac{\mathrm{d}}{\mathrm{d}t}(\hat{\mathbf{I}}\boldsymbol{v})
  =\hat{\mathbf{I}}\boldsymbol{a}+(\boldsymbol{v}\times^{*}\hat{\mathbf{I}}-\hat{\mathbf{I}}\,\boldsymbol{v}\times)\boldsymbol{v}
  =\hat{\mathbf{I}}\boldsymbol{a}+\boldsymbol{v}\times^{*}\hat{\mathbf{I}}\boldsymbol{v}.
  \label{eq:qk-spatial-eom}
\end{equation}
牛顿第二定律与欧拉方程被这\emph{一个}方程统一——这正是 \cref{ins:qk-why-spatial} “合一”动机的最终兑现。

\begin{note}[陀螺项的跨章对照]\label{note:qk-ocr-eom}
\cref{eq:qk-spatial-eom} 的末项 $\boldsymbol{v}\times^{*}\hat{\mathbf{I}}\boldsymbol{v}$ 即\emph{陀螺项}（空间动量叉乘）。\cref{ch:quad_mpc} 的单刚体凸 MPC \emph{忽略}这一陀螺项（小角速度下为二阶小量），把转动方程线性化为 $\dot{\boldsymbol{L}}\approx\mathbf{I}\dot{\boldsymbol{\omega}}$；而本章保留它以得到精确的全阶动力学。同一项，两种取舍，对照见 \cref{ins:qk-eom-contrast}。
\end{note}

\begin{insight}[要点：精确全阶（本章）vs 凸化简化（MPC 章）]\label{ins:qk-eom-contrast}
本章的 \cref{eq:qk-spatial-eom} 保留了完整陀螺项 $\boldsymbol{v}\times^{*}\hat{\mathbf{I}}\boldsymbol{v}$，是\emph{精确}的全阶刚体运动方程，用于 WBC 的逆动力学与高保真仿真；\cref{ch:quad_mpc} 为了让 MPC 每步退化成可实时求解的凸 QP，\emph{主动}丢弃陀螺项。理解“同一物理项在不同控制层被不同对待”，是把本章（运动学 / 全阶动力学地基）与 MPC 章（简化在线优化）、WBC 章（全阶任务分配）三者串成一条逻辑链的关键：\textbf{越靠下层（WBC）越追求精确，越靠上层在线优化（MPC）越追求凸性与速度}。
\end{insight}

\begin{note}[六维空间向量的常用性质（递推算法的代数底座）]\label{note:qk-spatial-props}
Plücker 定义虽繁，一旦接受，下列八条性质使整机递推（\cref{alg:qk-recursive-fk}）与组合刚体、力递推变得简洁，逐条标注其用处\cite{yu2021quadruped}：
\begin{enumerate}
  \item \textbf{相对速度}：刚体 $2$ 相对刚体 $1$ 的空间速度 $\boldsymbol{v}_{rel}=\boldsymbol{v}_2-\boldsymbol{v}_1$（加速度同样简单相减）——这是关节速度 $\boldsymbol{v}_{J,i}=\boldsymbol{v}_i-\boldsymbol{v}_{\lambda(i)}$（\cref{eq:qk-joint-vel}）的依据。
  \item \textbf{刚性连接}：两刚体若刚性连接，则空间速度\emph{相同}（一个量、无需分解平动转动）。
  \item \textbf{空间惯量可加}：刚性连接的两刚体组合惯量 $=\hat{\mathbf{I}}_1+\hat{\mathbf{I}}_2$，\emph{与相对位置无关}——组合刚体法（CRBA，\cref{ch:quad_wbc}）求质量矩阵 $\mathbf{M}$ 的基石。
  \item \textbf{空间力可加}：作用于同一刚体的两空间力等效于 $\hat{\boldsymbol{f}}_1+\hat{\boldsymbol{f}}_2$。
  \item \textbf{作用力 / 反作用力}：刚体 $1$ 对 $2$ 施空间力 $\hat{\boldsymbol{f}}$，则 $2$ 对 $1$ 施 $-\hat{\boldsymbol{f}}$——力递推（RNEA）里子刚体反作用项的来源。
  \item \textbf{微分封闭}：运动向量的微分仍是运动向量、力向量的微分仍是力向量（故 $\boldsymbol{v}\times$ 映 $\mathrm{M}^6\!\to\!\mathrm{M}^6$、$\boldsymbol{v}\times^{*}$ 映 $\mathrm{F}^6\!\to\!\mathrm{F}^6$）。
  \item \textbf{功率}：空间力 $\hat{\boldsymbol{f}}$ 作用于速度 $\boldsymbol{v}$ 的刚体，功率 $=\hat{\boldsymbol{f}}\cdot\boldsymbol{v}=\boldsymbol{v}^\top\hat{\boldsymbol{f}}$——\cref{deriv:qk-plucker-dual} 保功率推 $\mathbf{X}^{*}=\mathbf{X}^{-\top}$ 即用此。
  \item \textbf{动能}：空间惯量 $\hat{\mathbf{I}}$、空间速度 $\boldsymbol{v}$ 的刚体动能 $=\tfrac12\boldsymbol{v}^\top\hat{\mathbf{I}}\boldsymbol{v}$——CRBA 由“系统总动能 $=\sum_k\tfrac12\boldsymbol{v}_k^\top\hat{\mathbf{I}}_k\boldsymbol{v}_k=\tfrac12\dot{\boldsymbol{q}}^\top\mathbf{M}\dot{\boldsymbol{q}}$”比对出 $\mathbf{M}$。
\end{enumerate}
\end{note}

\subsection{关节模型与递推正运动学算法}
\label{subsec:qk-recursive-fk}

\paragraph{关节坐标变换与关节运动子空间。}
刚体 $i$ 相对父节点 $\lambda(i)$ 的变换分解为关节变换 $\mathbf{X}_J(i)$ 与固定的树变换 $\mathbf{X}_T(i)$，关节速度由关节运动子空间 $\mathbf{S}_i$ 映射\cite{yu2021quadruped}：
\begin{equation}
  {}^{i}\mathbf{X}_{\lambda(i)}=\mathbf{X}_J(i)\,\mathbf{X}_T(i),
  \qquad
  \boldsymbol{v}_{J,i}=\mathbf{S}_i\,\dot{\boldsymbol{q}}_i.
  \label{eq:qk-joint-model}
\end{equation}
这里 $\mathbf{X}_T(i)$ 是从父刚体系 $\mathcal{B}_{\lambda(i)}$ 到“关节零位重合系”$\mathcal{B}_{\lambda(i),i}$ 的\emph{常值}变换（由 \cref{exam:qk-dims} 的尺寸与坐标系位姿定出），$\mathbf{X}_J(i)$ 是随关节参数变化的关节变换。

\paragraph{关节速度与关节运动子空间 $\mathbf{S}_i$。}
关节速度定义为\emph{子刚体相对父刚体}的空间速度（即 \cref{note:qk-spatial-props} 性质 1）\cite{yu2021quadruped}：
\begin{equation}
  \boldsymbol{v}_{J,i}=\boldsymbol{v}_i-\boldsymbol{v}_{\lambda(i)}.
  \label{eq:qk-joint-vel}
\end{equation}
关节参数速度 $\dot{\boldsymbol{q}}_i$ 决定 $\boldsymbol{v}_{J,i}$，但二者维数\emph{不同}：$\boldsymbol{v}_{J,i}$ 恒为 $6$ 维空间速度，而 $\dot{\boldsymbol{q}}_i$ 的维数 $=$ 该关节自由度数（单自由度时为 $1$，浮动基座为 $6$），故 $\boldsymbol{v}_{J,i}\neq\dot{\boldsymbol{q}}_i$，须用映射 $\boldsymbol{v}_{J,i}=\mathbf{S}_i\dot{\boldsymbol{q}}_i$ 连接\cite{yu2021quadruped}。物理上：无约束时空间速度遍历整个 $\mathrm{M}^6$，受关节约束后 $\boldsymbol{v}_{J,i}$ 被限到 $\mathrm{M}^6$ 的一个\emph{子空间}（关节运动子空间），$\mathbf{S}_i$ 的每一列都是一个六维运动向量、张成该子空间，故称\textbf{关节运动子空间映射矩阵}。把 \cref{eq:qk-joint-vel} 代回即得速度递推 $\boldsymbol{v}_i=\boldsymbol{v}_{\lambda(i)}+\mathbf{S}_i\dot{\boldsymbol{q}}_i$（\cref{eq:qk-recursive} 速度行）。

浮动基座（虚拟 $6$-DoF 关节，\cref{ins:qk-floating-virtual}）的关节变换与子空间为\cite{yu2021quadruped}：
\begin{equation}
  \mathbf{X}_J(1)=\begin{bmatrix}{}^{\mathcal{B}}\mathbf{R}_{\mathcal{O}}&\mathbf{0}\\-{}^{\mathcal{B}}\mathbf{R}_{\mathcal{O}}\,({}^{\mathcal{O}}\boldsymbol{p}_{com})^{\wedge}&{}^{\mathcal{B}}\mathbf{R}_{\mathcal{O}}\end{bmatrix},
  \qquad {}^{1}\mathbf{S}_1=\mathbf{I}_{6\times6}.
  \label{eq:qk-floating-X}
\end{equation}
除浮动基座外，四足其余关节皆单自由度旋转关节：每腿第一关节绕 $x$、第二三关节绕 $y$。其关节变换为纯旋转（\cref{eq:qk-plucker-rot}），在子刚体系 $\mathcal{B}_i$ 下旋转轴即坐标轴，故子空间是对应的单位运动基\cite{yu2021quadruped}：
\begin{equation}
  \mathbf{X}_J(i)=\operatorname{diag}\big(\mathbf{R}_x^\top(q_i),\,\mathbf{R}_x^\top(q_i)\big),\ {}^{i}\mathbf{S}_i=[1\,0\,0\,0\,0\,0]^\top\ (\text{绕}\ x);
  \quad
  \mathbf{X}_J(i)=\operatorname{diag}\big(\mathbf{R}_y^\top(q_i),\,\mathbf{R}_y^\top(q_i)\big),\ {}^{i}\mathbf{S}_i=[0\,1\,0\,0\,0\,0]^\top\ (\text{绕}\ y).
  \label{eq:qk-single-dof}
\end{equation}
（绕 $z$ 同理，将 ${}^{i}\mathbf{S}_i$ 的第三元置 $1$。）注意 $\mathbf{X}_J$ 用 $\mathbf{R}^\top$ 而非 $\mathbf{R}$：变换矩阵把\emph{父系}量送到\emph{子系}，对应 ${}^{i}\mathbf{R}_{\lambda(i)}={}^{\lambda(i)}\mathbf{R}_i^\top$。

\paragraph{足底接触点变换。}
四条腿足底刚体编号 $\{\rho(1),\dots,\rho(4)\}=\{4,7,10,13\}$，足底接触点相对足端连杆系的（纯平移）变换为\cite{yu2021quadruped}：
\begin{equation}
  \mathbf{X}_C(i)=\begin{bmatrix}\mathbf{I}_3&\mathbf{0}\\-\big([0,0,-l_3]^\top\big)^{\wedge}&\mathbf{I}_3\end{bmatrix}.
  \label{eq:qk-foot-X}
\end{equation}

\paragraph{递推正运动学。}
沿运动学树从根（机身）向叶（足端）递推，逐刚体计算坐标变换、速度、加速度\cite{yu2021quadruped}：
\begin{equation}
  \boldsymbol{v}_i=\boldsymbol{v}_{\lambda(i)}+\mathbf{S}_i\dot{\boldsymbol{q}}_i,
  \qquad
  \boldsymbol{a}_i=\boldsymbol{a}_{\lambda(i)}+\mathbf{S}_i\ddot{\boldsymbol{q}}_i+\boldsymbol{v}_i\times\mathbf{S}_i\dot{\boldsymbol{q}}_i.
  \label{eq:qk-recursive}
\end{equation}
其中速度递推 $\boldsymbol{v}_i=\boldsymbol{v}_{\lambda(i)}+\boldsymbol{v}_{J,i}$ 表“子刚体速度 $=$ 父刚体速度 $+$ 关节相对速度”；加速度多出的 $\boldsymbol{v}_i\times\mathbf{S}_i\dot{\boldsymbol{q}}_i$ 是空间叉乘带来的科氏 / 离心耦合项。整理为算法：

\begin{algo}[运动学树递推正运动学（Featherstone）]\label{alg:qk-recursive-fk}
\textbf{初始化}：$\boldsymbol{v}_0=\mathbf{0}$，$\boldsymbol{a}_0=\mathbf{0}$，${}^{0}\mathbf{X}_0=\mathbf{I}_6$。\\
\textbf{对} $i=1,2,\dots,N$ 依树序（父先于子）执行：
\begin{enumerate}
  \item 关节变换：${}^{i}\mathbf{X}_{\lambda(i)}=\mathbf{X}_J(i)\,\mathbf{X}_T(i)$；
  \item 累积到根：${}^{i}\mathbf{X}_0={}^{i}\mathbf{X}_{\lambda(i)}\,{}^{\lambda(i)}\mathbf{X}_0$；
  \item 关节速度：$\boldsymbol{v}_{J,i}={}^{i}\mathbf{S}_i\,\dot{\boldsymbol{q}}_i$；
  \item 刚体速度：${}^{i}\boldsymbol{v}_i={}^{i}\mathbf{X}_{\lambda(i)}\,{}^{\lambda(i)}\boldsymbol{v}_{\lambda(i)}+\boldsymbol{v}_{J,i}$；
  \item 刚体加速度：${}^{i}\boldsymbol{a}_i={}^{i}\mathbf{X}_{\lambda(i)}\,{}^{\lambda(i)}\boldsymbol{a}_{\lambda(i)}+{}^{i}\mathbf{S}_i\,\ddot{\boldsymbol{q}}_i+{}^{i}\boldsymbol{v}_i\times{}^{i}\mathbf{S}_i\,\dot{\boldsymbol{q}}_i$。
\end{enumerate}
\end{algo}

\noindent{}该算法对浮动基座（$i=1$，$\mathbf{S}_1=\mathbf{I}_6$）与各腿单自由度关节\emph{一视同仁}，正是 \cref{ins:qk-floating-virtual} “浮动基座 $=$ 虚拟关节”带来的统一性红利。复杂度 $O(N)$，可扩展到任意刚体数。

% ════════════════════════════════════════════════════════════════
\section{全身控制雅可比}
\label{sec:qk-wbjac}

把递推 FK 推到“足端速度对全部关节速度的线性映射”，即得\emph{全身雅可比}——这是 \cref{ch:quad_wbc} 多任务 WBC 直接消费的对象。WBC 各级任务的雅可比 $\mathbf{J}_1,\mathbf{J}_4$ 及各级 $\dot{\mathbf{J}}_i\dot{\boldsymbol{q}}$ 项在那里被视为已知，本节在 Plücker 框架下把它们求出\cite{yu2021quadruped}。

\paragraph{运动学树的辅助节点集。}
除父节点集 $\lambda$（\cref{eq:qk-lambda}）外，递推算法还要三个集合\cite{yu2021quadruped}：$\kappa(i)$ 为“除 $0$ 号外、所有\emph{支撑起}刚体 $i$ 的关节（节点）集”（即根到 $i$ 路径上的关节）；$\mu(i)$ 为“以 $i$ 为\emph{父节点}的子节点集”（$i$ 的直接孩子）；$\nu(i)$ 为“以 $i$ 为\emph{根}的子树上所有节点集”。三者各有专责：$\kappa$ 装刚体雅可比（向根回溯）、$\mu$ 做力 / 惯量的向根累加（\cref{note:qk-spatial-props} 性质 3/5）、$\nu$ 定组合刚体范围。

\paragraph{刚体雅可比。}
对速度递推 \cref{eq:qk-recursive} 展开（沿 $\kappa(i)$ 累加），得刚体 $i$ 空间速度的\emph{非递推}式——它是到根路径上所有关节速度的线性叠加\cite{yu2021quadruped}：
\begin{equation}
  \boldsymbol{v}_i=\sum_{j\in\kappa(i)}\mathbf{S}_j\dot{\boldsymbol{q}}_j
  =\sum_{j=1}^{N}\epsilon_{ij}\mathbf{S}_j\dot{\boldsymbol{q}}_j
  =\underbrace{[\,\epsilon_{i1}\mathbf{S}_1\ \cdots\ \epsilon_{iN}\mathbf{S}_N\,]}_{{}_{b}\mathbf{J}_i}\dot{\boldsymbol{q}}
  \;=\;{}_{b}\mathbf{J}_i\,\dot{\boldsymbol{q}},
  \qquad
  \epsilon_{ij}=\begin{cases}1,&j\in\kappa(i)\\0,&\text{otherwise},\end{cases}
  \label{eq:qk-rigid-jacobian}
\end{equation}
其中 ${}_{b}\mathbf{J}_i$ 称\textbf{刚体雅可比}（关节空间广义速度 $\dot{\boldsymbol{q}}\in\mathbb{R}^{18}$ 到刚体空间速度 $\boldsymbol{v}_i\in\mathbb{R}^6$ 的 $6\times18$ 映射），指示量 $\epsilon_{ij}$ 自动“选出”路径上的关节列、其余列填零。具体在小腿刚体系 $\mathcal{B}_{\rho(i)}$ 下，从足端刚体沿 $\lambda$ 向根回溯、逐段累乘坐标变换并填列\cite{yu2021quadruped}：
\begin{algo}[小腿刚体雅可比 ${}_{b}\mathbf{J}_{\rho(i)}$ 的回溯装配]\label{alg:qk-rigid-jacobian}
\textbf{对}每条足 $i=1,\dots,N_C$：置 $j=\rho(i)$，${}^{\rho(i)}\mathbf{X}_j=\mathbf{I}_6$，${}_{b}\mathbf{J}_{\rho(i)}=\mathbf{0}$，列块 ${}_{b}\mathbf{J}_{\rho(i),\rho(i)}={}^{\rho(i)}\mathbf{S}_{\rho(i)}$；\\
\textbf{当} $\lambda(j)>0$：累乘 ${}^{\rho(i)}\mathbf{X}_{\lambda(j)}={}^{\rho(i)}\mathbf{X}_j\,{}^{j}\mathbf{X}_{\lambda(j)}$，令 $j=\lambda(j)$，填列块 ${}_{b}\mathbf{J}_{\rho(i),j}={}^{\rho(i)}\mathbf{X}_j\,{}^{j}\mathbf{S}_j$。
\end{algo}
\noindent{}列块 ${}_{b}\mathbf{J}_{\rho(i),j}$ 行数 $6$、列数 $=$ 第 $j$ 关节自由度数（机身列块 $6$ 列、各腿关节列块 $1$ 列），故拼成 $6\times18$。

\paragraph{足底系投影与任务雅可比。}
与定向本体系类似，在第 $i$ 足底建\emph{定向足底系} $\mathcal{Q}_i$：原点与足底系重合、\emph{方向同世界系}。设 ${}^{\rho(i)}\mathbf{R}_0$ 为固定基座系到第 $i$ 小腿刚体系的旋转（取自递推 FK 算出的 ${}^{\rho(i)}\mathbf{X}_0$ 右下 $3\times3$ 块），构建足底系到 $\mathcal{Q}_i$ 的变换并把刚体雅可比经足底接触点变换 $\mathbf{X}_C$、足底系投影 $\mathbf{X}_Q$ 送入 $\mathcal{Q}_i$\cite{yu2021quadruped}：
\begin{equation}
  \mathbf{X}_Q(i)=\operatorname{diag}\big({}^{\rho(i)}\mathbf{R}_0^\top,\,{}^{\rho(i)}\mathbf{R}_0^\top\big),
  \qquad
  {}_{b}^{\mathcal{Q}}\mathbf{J}_i=\mathbf{X}_Q(i)\,\mathbf{X}_C(i)\,{}_{b}\mathbf{J}_{\rho(i)}.
  \label{eq:qk-task-jacobian}
\end{equation}
因 $\mathcal{Q}_i$ 方向同世界、原点在足底接触点 $C_i$，${}_{b}^{\mathcal{Q}}\mathbf{J}_i$（$6\times18$）的\emph{后 $3$ 行}恰是 $C_i$ 在世界系下的\emph{平动}速度映射。WBC 第 $1$、$4$ 任务只控足底平动，故取后 $3$ 行 ${}_{b}^{\mathcal{Q}}\mathbf{J}_i'={}_{b}^{\mathcal{Q}}\mathbf{J}_i(4\!:\!6,\,1\!:\!18)$。再按支撑 / 摆动把各足 ${}_{b}^{\mathcal{Q}}\mathbf{J}_i'$ \emph{按列堆叠}\cite{yu2021quadruped}：
\begin{equation}
  \mathbf{J}_1=\operatorname*{stack}_{i:\ \text{支撑足}}{}_{b}^{\mathcal{Q}}\mathbf{J}_i',
  \qquad
  \mathbf{J}_4=\operatorname*{stack}_{i:\ \text{摆动足}}{}_{b}^{\mathcal{Q}}\mathbf{J}_i',
  \label{eq:qk-J1J4-assemble}
\end{equation}
即得 WBC 第 $1$（支撑足接触）、第 $4$（摆动足跟踪）任务雅可比；机身位置、姿态任务雅可比 $\mathbf{J}_2,\mathbf{J}_3$ 直接由浮动基座那 $6$ 列读出（常值）\cite{yu2021quadruped}。

\paragraph{$\dot{\mathbf{J}}\dot{\boldsymbol{q}}$ 项（偏置加速度）。}
对 $\boldsymbol{x}=\boldsymbol{f}(\boldsymbol{q})$、$\dot{\boldsymbol{x}}=\mathbf{J}\dot{\boldsymbol{q}}$ 再求一次时间导，得任务空间加速度 $\ddot{\boldsymbol{x}}=\mathbf{J}\ddot{\boldsymbol{q}}+\dot{\mathbf{J}}\dot{\boldsymbol{q}}$；偏置项 $\dot{\mathbf{J}}\dot{\boldsymbol{q}}$ 即\emph{令 $\ddot{\boldsymbol{q}}=0$} 时的任务空间加速度。据 \cref{eq:qk-recursive} 加速度行去掉 $\mathbf{S}_i\ddot{\boldsymbol{q}}_i$ 项，得各刚体的“速度积加速度”（velocity-product，VP）递推\cite{yu2021quadruped}：
\begin{equation}
  \boldsymbol{a}_i^{vp}=\boldsymbol{a}_{\lambda(i)}^{vp}+\boldsymbol{v}_i\times\mathbf{S}_i\dot{\boldsymbol{q}}_i,
  \qquad \boldsymbol{a}_1^{vp}=\boldsymbol{a}_0^{vp}+\boldsymbol{v}_1\times\mathbf{S}_1\dot{\boldsymbol{q}}_1=\boldsymbol{v}_1\times\boldsymbol{v}_1=\mathbf{0},
  \label{eq:qk-Jdot-q}
\end{equation}
浮动基座项为零是因 $\mathbf{S}_1\dot{\boldsymbol{q}}_1=\boldsymbol{v}_1$ 而 $\boldsymbol{v}_1\times\boldsymbol{v}_1=\mathbf{0}$。由此立得\textbf{机身位置 / 姿态任务（任务 $2,3$）的偏置项恒零}\cite{yu2021quadruped}：
\begin{equation}
  \dot{\mathbf{J}}_2\dot{\boldsymbol{q}}=\mathbf{0}_{3\times1},
  \qquad
  \dot{\mathbf{J}}_3\dot{\boldsymbol{q}}=\mathbf{0}_{3\times1}.
  \label{eq:qk-J23-zero}
\end{equation}
对足底任务（$1,4$），把第 $i$ 小腿的 ${}^{i}\boldsymbol{a}_{\rho(i)}^{vp}$ 与 ${}^{\rho(i)}\boldsymbol{v}_{\rho(i)}$ 经 $\mathbf{X}_Q\mathbf{X}_C$ 转到 $\mathcal{Q}_i$，取平动分量 $\boldsymbol{a}_{\rho(i)}^{C,i},\boldsymbol{v}_{\rho(i)}^{C,i}$ 与转动分量 $\boldsymbol{\omega}_{\rho(i)}$，由空间加速度定义（\cref{eq:qk-spatial-acc}）得足底接触点世界系偏置加速度\cite{yu2021quadruped}：
\begin{equation}
  {}^{\mathcal{O}}\boldsymbol{a}_{C,i}=\boldsymbol{a}_{\rho(i)}^{C,i}+\boldsymbol{\omega}_{\rho(i)}\times\boldsymbol{v}_{\rho(i)}^{C,i},
  \label{eq:qk-foot-bias-acc}
\end{equation}
再按支撑 / 摆动堆叠成 $\dot{\mathbf{J}}_1\dot{\boldsymbol{q}}$、$\dot{\mathbf{J}}_4\dot{\boldsymbol{q}}$（与 \cref{eq:qk-J1J4-assemble} 同结构）。这 $6$ 个矩阵正是 \cref{ch:quad_wbc} WBC 各级任务所缺的最后拼图。

\paragraph{力侧对偶：$\boldsymbol{\tau}=\mathbf{S}^\top\hat{\boldsymbol{f}}_J$（虚功原理）。}
单腿静力学的 $\boldsymbol{\tau}=\mathbf{J}^\top\mathbf{F}$（\cref{thm:qk-duality}）在 Plücker 框架下有一个更基本的对偶：设父刚体经关节作用到子刚体的空间力为 $\hat{\boldsymbol{f}}_J$，关节约束力在允许运动方向不做功；给关节虚位移 $\delta\boldsymbol{q}$，则 $\boldsymbol{v}_J$ 方向虚位移 $\delta\boldsymbol{d}=\mathbf{S}\,\delta\boldsymbol{q}$，虚功原理对任意 $\delta\boldsymbol{q}$ 给出\cite{yu2021quadruped}：
\begin{equation}
  \boldsymbol{\tau}^\top\delta\boldsymbol{q}=\hat{\boldsymbol{f}}_J^\top\delta\boldsymbol{d}=\hat{\boldsymbol{f}}_J^\top\mathbf{S}\,\delta\boldsymbol{q}
  \;\Longrightarrow\;
  \boldsymbol{\tau}=\mathbf{S}^\top\hat{\boldsymbol{f}}_J.
  \label{eq:qk-virtual-work}
\end{equation}
这与 \cref{thm:qk-duality} 同出一辙（功率 / 虚功守恒 $\Rightarrow$ 转置映射），但把对象从“足端力 / 足端雅可比”推广到“任意关节力 / 关节运动子空间 $\mathbf{S}$”——递推牛顿--欧拉逆动力学（\cref{ch:quad_wbc}）由它从空间力反算关节力矩。这是 \cref{ins:qk-duality}（K6）的全身版本。

\paragraph{浮动基座动力学标准形（衔接预告）。}
以上全身雅可比将进入浮动基座动力学标准形\cite{yu2021quadruped}：
\begin{equation}
  \mathbf{M}(\boldsymbol{q})\ddot{\boldsymbol{q}}+\mathbf{C}(\boldsymbol{q},\dot{\boldsymbol{q}})=\mathbf{S}_\tau\boldsymbol{\tau}+\mathbf{J}_c^\top\,{}^{\mathcal{O}}\mathbf{f}_c,
  \label{eq:qk-floating-dynamics}
\end{equation}
其中 $\mathbf{M}$ 为质量矩阵、$\mathbf{C}$ 含科氏 / 离心 / 重力、$\mathbf{S}_\tau$ 为驱动选择矩阵（浮动基座 $6$ 维不可驱）、$\mathbf{J}_c^\top{}^{\mathcal{O}}\mathbf{f}_c$ 为接触力。本式的完整动力学推导与求解留给 \cref{ch:quad_wbc}，本章只引出接口。

% ════════════════════════════════════════════════════════════════
\section{奇异性与零空间理论：WBC 的数学地基}
\label{sec:qk-nullspace}

本节是本章理论高潮，把“冗余 / 零空间 / 伪逆 / 阻尼”连同完整证明讲透，直接作为 \cref{ch:quad_wbc} 多任务全身控制的数学地基。

\subsection{冗余与零空间}
\label{subsec:qk-redundancy}

一阶微分运动学方程与雅可比的零空间定义为\cite{yu2021quadruped}：
\begin{equation}
  \dot{\boldsymbol{x}}=\mathbf{J}\dot{\boldsymbol{q}},
  \qquad
  \mathcal{N}(\mathbf{J})=\{\boldsymbol{w}\in\mathbb{R}^N\mid \mathbf{J}\boldsymbol{w}=\mathbf{0}\}.
  \label{eq:qk-null-def}
\end{equation}
当关节数 $N$ 大于任务维 $M$（如 $7$-DoF 臂够 $6$-DoF 目标），系统\emph{冗余}，$\mathcal{N}(\mathbf{J})$ 非平凡（含非零向量）。

\begin{insight}[K7：零空间 $=$ 冗余的“免费内部运动”（WBC 多任务本质）]\label{ins:qk-nullspace-free}
零空间的物理意义是理解 WBC 的钥匙：满足 $\mathbf{J}\dot{\boldsymbol{q}}_0=\mathbf{0}$ 的关节速度 $\dot{\boldsymbol{q}}_0\in\mathcal{N}(\mathbf{J})$ 是\textbf{“关节在动、末端却不动”的内部运动}——它\emph{不改变}主任务（末端位姿），却能让机器人调整自身构型。所以零空间运动是\emph{免费}的：在完美执行主任务的前提下，可以“顺便”用它去做次要任务（避障、避关节限位、优化操作度）。\textbf{这正是 WBC 多任务优先级的本质}——主任务占满雅可比的行空间，次任务被投影到零空间里执行，互不干扰。本章把它讲透，就是为 \cref{ch:quad_wbc} 的“任务优先级 / 零空间投影”铺好数学地基。
\end{insight}

\subsection{通解结构与零空间投影矩阵}
\label{subsec:qk-general-solution}

冗余系统 $\dot{\boldsymbol{x}}=\mathbf{J}\dot{\boldsymbol{q}}$ 的解不唯一，其通解由\emph{特解}加\emph{齐次解}构成\cite{yu2021quadruped}：
\begin{equation}
  \dot{\boldsymbol{q}}=\underbrace{\mathbf{J}^{+}\dot{\boldsymbol{x}}}_{\dot{\boldsymbol{q}}_p\ (\text{特解})}+\underbrace{\mathbf{N}\dot{\boldsymbol{q}}_{\forall}}_{\dot{\boldsymbol{q}}_0\in\mathcal{N}(\mathbf{J})\ (\text{内部运动})},
  \qquad
  \mathbf{J}^{+}=\mathbf{J}^\top(\mathbf{J}\mathbf{J}^\top)^{-1},
  \quad
  \mathbf{N}=\mathbf{I}-\mathbf{J}^{+}\mathbf{J}.
  \label{eq:qk-general-solution}
\end{equation}
其中 $\mathbf{J}^{+}$ 是（右）Moore--Penrose 伪逆（满足 $\mathbf{J}\mathbf{J}^{+}=\mathbf{I}_M$），$\dot{\boldsymbol{q}}_{\forall}$ 为任意向量、$\mathbf{N}$ 把它投影到零空间。下面先说明 $\mathbf{N}$ 的形式\emph{从哪来}，再证它是正交投影。

\paragraph{右逆与 $\mathbf{N}$ 形式的来历。}
冗余下 $\mathbf{J}$（$M\times N$，$M<N$）非方阵、无普通逆。但可定义\emph{右逆} $\mathbf{G}$（$N\times M$）满足 $\mathbf{J}\mathbf{G}=\mathbf{I}_M$，则 $\dot{\boldsymbol{q}}_p=\mathbf{G}\dot{\boldsymbol{x}}$ 即一个特解（验证：$\mathbf{J}\dot{\boldsymbol{q}}_p=\mathbf{J}\mathbf{G}\dot{\boldsymbol{x}}=\dot{\boldsymbol{x}}$）\cite{yu2021quadruped}。\emph{右逆有无穷多}，$\mathbf{J}^{+}=\mathbf{J}^\top(\mathbf{J}\mathbf{J}^\top)^{-1}$ 只是其中一个特例（数学上称 Moore--Penrose 伪逆，下节证其“能量最优”）。一旦有了任一右逆，零空间投影矩阵的形式即\emph{被它逼出}：把 $\mathbf{I}_M=\mathbf{J}\mathbf{J}^{+}$ 代入恒等变形
\begin{equation}
  \mathbf{0}=\mathbf{J}-\mathbf{J}=\mathbf{J}-\mathbf{I}_M\mathbf{J}=\mathbf{J}-(\mathbf{J}\mathbf{J}^{+})\mathbf{J}=\mathbf{J}(\mathbf{I}-\mathbf{J}^{+}\mathbf{J}),
  \label{eq:qk-N-from-rightinv}
\end{equation}
最右括号 $\mathbf{I}-\mathbf{J}^{+}\mathbf{J}$ 被 $\mathbf{J}$ 左乘归零——这恰是零空间投影矩阵的定义性质 $\mathbf{J}\mathbf{N}=\mathbf{0}$，故取 $\mathbf{N}=\mathbf{I}-\mathbf{J}^{+}\mathbf{J}$\cite{yu2021quadruped}。代入 \cref{eq:qk-general-solution} 通解，左乘 $\mathbf{J}$ 即验证其正确：$\mathbf{J}\dot{\boldsymbol{q}}=\mathbf{J}\mathbf{J}^{+}\dot{\boldsymbol{x}}+\mathbf{J}\mathbf{N}\dot{\boldsymbol{q}}_{\forall}=\dot{\boldsymbol{x}}+\mathbf{0}=\dot{\boldsymbol{x}}$，不论 $\dot{\boldsymbol{q}}_{\forall}$ 取何值，主任务 $\dot{\boldsymbol{x}}$ 始终满足——这就是 \cref{ins:qk-nullspace-free}“免费内部运动”在代数上的兑现。

\begin{theorem}[零空间投影矩阵 $\mathbf{N}=\mathbf{I}-\mathbf{J}^{+}\mathbf{J}$]\label{thm:qk-nullspace-proj}
矩阵 $\mathbf{N}=\mathbf{I}-\mathbf{J}^{+}\mathbf{J}$ 把任意 $\dot{\boldsymbol{q}}_{\forall}$ 映射到 $\mathcal{N}(\mathbf{J})$，即对任意 $\dot{\boldsymbol{q}}_{\forall}$ 有 $\mathbf{J}\mathbf{N}\dot{\boldsymbol{q}}_{\forall}=\mathbf{0}$；且 $\mathbf{N}$ 是到 $\mathcal{N}(\mathbf{J})$ 的\emph{正交}投影（$\mathbf{N}^2=\mathbf{N}$、$\mathbf{N}^\top=\mathbf{N}$）。
\end{theorem}
\begin{proof}
首先验证投影到零空间：用 $\mathbf{J}\mathbf{J}^{+}=\mathbf{J}\mathbf{J}^\top(\mathbf{J}\mathbf{J}^\top)^{-1}=\mathbf{I}_M$，
\begin{equation}
  \mathbf{J}\mathbf{N}=\mathbf{J}(\mathbf{I}-\mathbf{J}^{+}\mathbf{J})=\mathbf{J}-(\mathbf{J}\mathbf{J}^{+})\mathbf{J}=\mathbf{J}-\mathbf{J}=\mathbf{0},
  \label{eq:qk-JN-zero}
\end{equation}
故 $\mathbf{J}(\mathbf{N}\dot{\boldsymbol{q}}_{\forall})=\mathbf{0}$，即 $\mathbf{N}\dot{\boldsymbol{q}}_{\forall}\in\mathcal{N}(\mathbf{J})$。其次幂等：
\[
  \mathbf{N}^2=(\mathbf{I}-\mathbf{J}^{+}\mathbf{J})(\mathbf{I}-\mathbf{J}^{+}\mathbf{J})
  =\mathbf{I}-2\mathbf{J}^{+}\mathbf{J}+\mathbf{J}^{+}(\mathbf{J}\mathbf{J}^{+})\mathbf{J}
  =\mathbf{I}-2\mathbf{J}^{+}\mathbf{J}+\mathbf{J}^{+}\mathbf{J}=\mathbf{N},
\]
用了 $\mathbf{J}\mathbf{J}^{+}=\mathbf{I}_M$。最后对称：$\mathbf{J}^{+}\mathbf{J}=\mathbf{J}^\top(\mathbf{J}\mathbf{J}^\top)^{-1}\mathbf{J}$，而 $(\mathbf{J}\mathbf{J}^\top)^{-1}$ 对称、$\mathbf{J}^\top(\cdots)\mathbf{J}$ 形如 $\mathbf{A}^\top\mathbf{S}\mathbf{A}$ 故对称，从而 $\mathbf{N}^\top=\mathbf{N}$。幂等 $+$ 对称 $\Rightarrow$ 正交投影。\qed
\end{proof}

\begin{insight}[要点：投影的“最优近似”含义]\label{ins:qk-proj-best}
$\mathbf{N}$ 是正交投影意味着 $\mathbf{N}\dot{\boldsymbol{q}}_{\forall}$ 是 $\dot{\boldsymbol{q}}_{\forall}$ 在零空间里\emph{距离最近}的向量（投影误差与零空间正交）。由此可把零空间投影解读为“\emph{在完美执行主任务的前提下，尽力近似次要任务期望}”：次任务的期望关节速度 $\dot{\boldsymbol{q}}_{\forall}$ 未必落在零空间，$\mathbf{N}$ 取其零空间分量——这正是 WBC 严格优先级（strict hierarchy）的几何实现。
\end{insight}

\begin{note}[零空间投影即“过滤器”：为何它是\emph{唯一}的正交投影]\label{note:qk-proj-filter}
下面一条更细的逻辑链解释 $\mathbf{N}$ 何以是\emph{唯一}选择\cite{yu2021quadruped}：引入次要任务（如避障）期望的任意关节速度 $\dot{\boldsymbol{q}}_{\forall}$，不能直接执行——其中某些分量会经 $\mathbf{J}$ 产生末端速度、\emph{破坏}高优先级主任务。于是需要一个“\textbf{过滤器}”$\mathbf{N}$，把 $\dot{\boldsymbol{q}}_{\forall}$ 投到零空间得 $\dot{\boldsymbol{q}}_0=\mathbf{N}\dot{\boldsymbol{q}}_{\forall}$，使 $\mathbf{J}\dot{\boldsymbol{q}}_0=\mathbf{0}$。\textbf{要求一}：无论 $\dot{\boldsymbol{q}}_{\forall}$ 取何值都须落零空间，故 $\mathbf{J}\mathbf{N}=\mathbf{0}$——这逼出 $\mathbf{N}$ 必是零空间矩阵（\cref{eq:qk-N-from-rightinv}）。\textbf{要求二}：过滤后“尽量保留原意图”，即投影后向量与原向量\emph{欧氏距离最小}——这正是\emph{正交投影}（误差 $\dot{\boldsymbol{q}}_{\forall}-\dot{\boldsymbol{q}}_0\perp\mathcal{N}(\mathbf{J})$）。给定子空间，正交投影矩阵\emph{唯一}（由 $\mathbf{N}^2=\mathbf{N},\mathbf{N}^\top=\mathbf{N}$ 与值域 $=\mathcal{N}(\mathbf{J})$ 唯一确定，\cref{thm:qk-nullspace-proj} 已证 $\mathbf{I}-\mathbf{J}^{+}\mathbf{J}$ 满足），故 $\mathbf{N}=\mathbf{I}-\mathbf{J}^{+}\mathbf{J}$ 是\emph{唯一}的零空间正交投影。\textbf{要点}：投影的意义\emph{不是}完美执行次任务，而是“\emph{在完美执行主任务的前提下，尽最大努力近似次任务}”——\cref{ch:quad_wbc} 的严格任务优先级即由“逐级把低优先任务投到高优先任务零空间”实现。
\end{note}

\subsection{Moore--Penrose 伪逆 \texorpdfstring{$=$}{=} 最小范数解}
\label{subsec:qk-mp-pinv}

\begin{insight}[K8：MP 伪逆 $=$ 能量最优（全身控制真正基石）]\label{ins:qk-mp-energy}
冗余系统有无穷多解满足 $\mathbf{J}\dot{\boldsymbol{q}}=\dot{\boldsymbol{x}}$。在这无穷多右逆中，为什么独选 Moore--Penrose 伪逆 $\mathbf{J}^{+}=\mathbf{J}^\top(\mathbf{J}\mathbf{J}^\top)^{-1}$？物理答案是：\textbf{MP 伪逆给出的是\emph{最小关节速度模长}的解}，物理上对应最节能、最平滑的运动。这一“能量最优”性质由拉格朗日乘子法严格导出（下证），使 MP 伪逆成为全身控制的\emph{真正基石}。
\end{insight}

\begin{theorem}[MP 伪逆的最小范数性质]\label{thm:qk-mp-minnorm}
在约束 $\mathbf{J}\dot{\boldsymbol{q}}=\dot{\boldsymbol{x}}$（$\mathbf{J}$ 行满秩）下，最小化 $\tfrac12\|\dot{\boldsymbol{q}}\|^2$ 的唯一解为 $\dot{\boldsymbol{q}}=\mathbf{J}^\top(\mathbf{J}\mathbf{J}^\top)^{-1}\dot{\boldsymbol{x}}=\mathbf{J}^{+}\dot{\boldsymbol{x}}$。
\end{theorem}
\begin{proof}
构造拉格朗日函数（乘子 $\boldsymbol{\lambda}\in\mathbb{R}^M$）\cite{yu2021quadruped}：
\begin{equation}
  \mathcal{L}(\dot{\boldsymbol{q}},\boldsymbol{\lambda})=\tfrac12\dot{\boldsymbol{q}}^\top\dot{\boldsymbol{q}}-\boldsymbol{\lambda}^\top(\mathbf{J}\dot{\boldsymbol{q}}-\dot{\boldsymbol{x}}).
  \label{eq:qk-lagrangian}
\end{equation}
平稳条件（对 $\dot{\boldsymbol{q}}$ 与 $\boldsymbol{\lambda}$ 求梯度置零）：
\begin{equation}
  \nabla_{\dot{\boldsymbol{q}}}\mathcal{L}=\dot{\boldsymbol{q}}-\mathbf{J}^\top\boldsymbol{\lambda}=\mathbf{0}\ \Rightarrow\ \dot{\boldsymbol{q}}=\mathbf{J}^\top\boldsymbol{\lambda};
  \qquad
  \nabla_{\boldsymbol{\lambda}}\mathcal{L}=-(\mathbf{J}\dot{\boldsymbol{q}}-\dot{\boldsymbol{x}})=\mathbf{0}\ \Rightarrow\ \mathbf{J}\dot{\boldsymbol{q}}=\dot{\boldsymbol{x}}.
  \label{eq:qk-stationary}
\end{equation}
将 $\dot{\boldsymbol{q}}=\mathbf{J}^\top\boldsymbol{\lambda}$ 代入约束 $\mathbf{J}\dot{\boldsymbol{q}}=\dot{\boldsymbol{x}}$：
\begin{equation}
  \mathbf{J}\mathbf{J}^\top\boldsymbol{\lambda}=\dot{\boldsymbol{x}}
  \ \Rightarrow\ \boldsymbol{\lambda}=(\mathbf{J}\mathbf{J}^\top)^{-1}\dot{\boldsymbol{x}}
  \quad(\mathbf{J}\text{ 行满秩 }\Rightarrow\mathbf{J}\mathbf{J}^\top\text{ 可逆}).
  \label{eq:qk-lambda-solve}
\end{equation}
回代得
\begin{equation}
  \dot{\boldsymbol{q}}=\mathbf{J}^\top(\mathbf{J}\mathbf{J}^\top)^{-1}\dot{\boldsymbol{x}}=\mathbf{J}^{+}\dot{\boldsymbol{x}}.
  \label{eq:qk-mp-result}
\end{equation}
目标函数为严格凸（海森 $=\mathbf{I}\succ0$）、约束为线性，故平稳点即\emph{唯一}全局最小。该解位于 $\mathbf{J}$ 的行空间（$\dot{\boldsymbol{q}}=\mathbf{J}^\top\boldsymbol{\lambda}$），与零空间正交——即 \cref{eq:qk-general-solution} 通解中 $\dot{\boldsymbol{q}}_0=\mathbf{0}$ 的那一支。\qed
\end{proof}

\paragraph{秩--零度定理（冗余维数）。}
输入空间维数等于值域维加零空间维（\cref{ch:matderiv}）\cite{yu2021quadruped}：
\begin{equation}
  N=\operatorname{rank}(\mathbf{J})+\dim\mathcal{N}(\mathbf{J}).
  \label{eq:qk-rank-nullity}
\end{equation}
非奇异（$\operatorname{rank}\mathbf{J}=M$）时 $\dim\mathcal{N}(\mathbf{J})=N-M>0$，故\emph{必}存在可用的内部运动——这从维数上保证了 \cref{ins:qk-nullspace-free} 的“免费运动”不是空谈。

\subsection{奇异性与阻尼最小二乘}
\label{subsec:qk-singularity}

\begin{insight}[K9：近奇异须阻尼最小二乘（数值稳健化）]\label{ins:qk-damped}
\cref{thm:qk-mp-minnorm} 的 $\mathbf{J}^{+}=\mathbf{J}^\top(\mathbf{J}\mathbf{J}^\top)^{-1}$ 要求 $\mathbf{J}\mathbf{J}^\top$ 可逆。但在\emph{近奇异}构型下（如腿伸直、两关节轴共线），$\mathbf{J}\mathbf{J}^\top$ 虽理论可逆，其行列式 $\det(\mathbf{J}\mathbf{J}^\top)\to0$ 成\emph{病态}——求逆数值溢出，解出的关节速度 $\dot{\boldsymbol{q}}$ 爆炸（微小的任务速度要求无穷大的关节速度）。工程标准对策是\textbf{阻尼最小二乘}（damped least squares，又称 Levenberg--Marquardt 正则）\cite{yu2021quadruped}：
\begin{equation}
  \mathbf{J}^{+}=\mathbf{J}^\top\big(\mathbf{J}\mathbf{J}^\top+\lambda^2\mathbf{I}\big)^{-1}.
  \label{eq:qk-damped-pinv}
\end{equation}
加入 $\lambda^2\mathbf{I}$（$\lambda>0$ 为阻尼因子）保证 $\mathbf{J}\mathbf{J}^\top+\lambda^2\mathbf{I}\succ0$ \emph{始终}可逆、特征值有下界 $\lambda^2$，从而求逆数值稳定。代价是在奇异附近\emph{牺牲少量跟踪精度}换取稳健（解不再精确满足 $\mathbf{J}\dot{\boldsymbol{q}}=\dot{\boldsymbol{x}}$）。\cref{eq:qk-damped-pinv} 等价于最小化 $\|\mathbf{J}\dot{\boldsymbol{q}}-\dot{\boldsymbol{x}}\|^2+\lambda^2\|\dot{\boldsymbol{q}}\|^2$ 的正则最小二乘——精度与稳健的权衡由 $\lambda$ 调节（\cref{ch:nlopt}）。这是本章开篇前置自测第 6 问的答案，也是欧拉角反面教材（\cref{ins:qk-euler-only-hmi}）与 WBC 章数值稳健性的共同基础。
\end{insight}

\begin{example}[雅可比奇异性的数值演示]\label{exam:qk-singular}
对单腿 $3\times3$ 雅可比 \cref{eq:qk-jacobian-def}，操作度 $w=\sqrt{\det(\mathbf{J}\mathbf{J}^\top)}=|\det\mathbf{J}|$ 度量离奇异的远近。远奇异构型（关节角适中、$\det\mathbf{J}$ 绝对值大）时直接求逆 \cref{eq:qk-jacobian-invert} 稳定；当腿接近\emph{完全伸直}（大腿小腿共线，$\theta_3\to0$），$\mathbf{J}$ 的两列趋于线性相关，$\det\mathbf{J}\to0$，$\mathbf{J}^{-1}$ 元素发散——此时用 \cref{eq:qk-damped-pinv} 阻尼伪逆，取 $\lambda$ 约为 $0.01\sim0.1$（量纲随 $\mathbf{J}$ 缩放），关节速度被限在合理范围。工程上常令 $\lambda$ 随 $w$ 减小而增大（仅在近奇异时介入），兼顾远奇异区的精度与近奇异区的稳健。
\end{example}

\begin{pitfall}[误区：以为 $\mathbf{J}\mathbf{J}^\top$ 可逆就万事大吉]
“理论可逆”不等于“数值可逆”。近奇异时 $\mathbf{J}\mathbf{J}^\top$ 的最小特征值趋零、条件数趋无穷，浮点求逆会放大舍入误差到不可用。判据应是\emph{条件数} / \emph{操作度} $w$ 而非“行列式是否恰为零”——后者几乎永远不为零，却早已病态。务必用阻尼伪逆 \cref{eq:qk-damped-pinv} 或基于 SVD 的截断伪逆兜底。
\end{pitfall}

% ════════════════════════════════════════════════════════════════
\section{四足整机运动学}
\label{sec:qk-wholebody}

把单腿运动学、雅可比与坐标变换整合，得到“机身位姿$\leftrightarrow$各足端”与“整机足端速度”的整机关系——这是步态控制（\cref{ch:quad_gait}）与状态估计（\cref{ch:quad_est}）的直接接口。

\subsection{机身姿态 \texorpdfstring{$\leftrightarrow$}{<->} 足端位置}
\label{subsec:qk-body-foot}

站立或支撑相中，足端在世界系\emph{不滑移}（位置固定为常量）。给定期望机身位姿 $\mathbf{T}_{sb}$，反解各足端在机身系的坐标，再走单腿 IK 即得关节角\cite{unitree2023quadruped}：
\begin{equation}
  \boldsymbol{p}_{si}=\boldsymbol{p}_{bi}(0)-\boldsymbol{p}_{b0}(0),
  \quad
  \mathbf{T}_{sb}=\begin{bmatrix}\mathbf{R}_z(\psi)\mathbf{R}_y(\theta)\mathbf{R}_x(\phi)&-\boldsymbol{p}_{b0}(0)+\boldsymbol{p}_d\\\mathbf{0}^\top&1\end{bmatrix},
  \quad
  \boldsymbol{p}_{bi}=\mathbf{T}_{sb}^{-1}\boldsymbol{p}_{si}.
  \label{eq:qk-body-foot}
\end{equation}
逻辑链：机身位姿变 $\Rightarrow$ 各足端在世界系 $\{s\}$ 仍不动（常量 $\boldsymbol{p}_{si}$）$\Rightarrow$ 经 $\mathbf{T}_{sb}^{-1}$ 反解机身系足端坐标 $\boldsymbol{p}_{bi}$ $\Rightarrow$ 单腿 IK（\cref{sec:qk-ik}）得关节角。这是“用关节控机身姿态”的直接落地。

\subsection{整机足端速度：牵连速度 \texorpdfstring{$+$}{+} 关节速度}
\label{subsec:qk-foot-vel}

\begin{insight}[K11：机身系是“瞬时重合的静止惯性系”（理解整机速度的钥匙）]\label{ins:qk-static-body}
一个微妙却关键的认知是：在推导足端速度时，机身系 $\{b\}$ \textbf{并\emph{不}随机身一起运动}，而是被取为\emph{当前时刻恰与机身位姿重合的、静止的惯性系}\cite{unitree2023quadruped}。为什么必须这样取？反证：若 $\{b\}$ 随机身运动（随动系），则关节静止时足端相对该系的速度\emph{恒为零}，无法定义“牵连速度”这一概念——而我们恰恰要用“机身运动牵连足端”加“关节运动驱动足端”两部分合成总速度。因此把 $\{b\}$ 取为瞬时重合的静止系，才能让足端速度自然分解为\emph{牵连速度 $+$ 关节速度}两项。这与 \cref{ch:quad_mpc} 取定向本体系 $\mathcal{P}$（方向同世界、原点随机身）的动机同源，是理解整机速度公式的钥匙。
\end{insight}

\paragraph{刚体上一点的速度。}
由 \cref{ins:qk-static-body} 取 $\{b\}$ 为瞬时静止惯性系，刚体上一点 $P$ 的速度 / 加速度为（机体角速度，右扰动主线）\cite{unitree2023quadruped}：
\begin{equation}
  \dot{\boldsymbol{p}}_P=\boldsymbol{v}_b+\boldsymbol{\omega}_b^{\wedge}\boldsymbol{p}_P,
  \qquad
  \ddot{\boldsymbol{p}}_P=\dot{\boldsymbol{v}}_b+\dot{\boldsymbol{\omega}}_b^{\wedge}\boldsymbol{p}_P+\boldsymbol{\omega}_b^{\wedge}\boldsymbol{\omega}_b^{\wedge}\boldsymbol{p}_P.
  \label{eq:qk-point-vel}
\end{equation}

\paragraph{足端总速度（牵连 $+$ 关节）。}
足端速度由“机身运动牵连项 $\boldsymbol{v}_b+\boldsymbol{\omega}_b^{\wedge}\boldsymbol{p}_{bf}$”与“关节驱动项 $\mathbf{J}\dot{\boldsymbol{\theta}}$”合成。先在机身系合成，再经 ${}^{\mathcal{O}}\mathbf{R}_{\mathcal{B}}=\mathbf{R}_{sb}$ 转世界系\cite{unitree2023quadruped}：
\begin{equation}
  \boldsymbol{v}_{be}=\boldsymbol{v}_b+\boldsymbol{\omega}_b^{\wedge}\boldsymbol{p}_{bf},
  \qquad
  \boldsymbol{v}_{bj}=\mathbf{J}\dot{\boldsymbol{\theta}},
  \label{eq:qk-foot-vel-body}
\end{equation}
\begin{equation}
  \boldsymbol{v}_{sf}=\boldsymbol{v}_s+\mathbf{R}_{sb}\big(\boldsymbol{\omega}_b^{\wedge}\boldsymbol{p}_{bf}+\mathbf{J}\dot{\boldsymbol{\theta}}\big),
  \qquad
  \boldsymbol{v}_{sfB}=\boldsymbol{v}_{sf}-\boldsymbol{v}_s=\mathbf{R}_{sb}\big(\boldsymbol{\omega}_b^{\wedge}\boldsymbol{p}_{bf}+\mathbf{J}\dot{\boldsymbol{\theta}}\big).
  \label{eq:qk-foot-vel-world}
\end{equation}

\begin{insight}[要点：足端相对机身速度是状态估计的可测量]\label{ins:qk-vsfB}
\cref{eq:qk-foot-vel-world} 中机身绝对速度 $\boldsymbol{v}_s$ 在真机上\emph{不可直接测量}（无外部定位时），但\emph{足端相对机身速度} $\boldsymbol{v}_{sfB}=\mathbf{R}_{sb}(\boldsymbol{\omega}_b^{\wedge}\boldsymbol{p}_{bf}+\mathbf{J}\dot{\boldsymbol{\theta}})$ \emph{完全由可测量算出}（机体角速度 $\boldsymbol{\omega}_b$ 来自 IMU、关节角 / 角速度来自编码器、$\mathbf{J}$ 与 $\boldsymbol{p}_{bf}$ 来自本章运动学）。这正是足式里程计 / 状态估计的核心观测：支撑足在世界系不动（$\boldsymbol{v}_{sf}\approx\mathbf{0}$），于是 $\boldsymbol{v}_s\approx-\boldsymbol{v}_{sfB}$ 反推机身速度。本章的运动学因此直接喂给 \cref{ch:quad_est} 的足式状态估计——这是“运动学是估计的前向模型”的体现。
\end{insight}

\subsection{电机零点与模型零点}
\label{subsec:qk-motor-zero}

\begin{practice}[K12：电机零点 $\neq$ 模型零点（实物与模型的连接）]\label{prac:qk-motor-zero}
本章所有 FK/IK 公式中的关节角 $q_i$（或 $\theta_i$）都是\emph{模型零点}下的角度（即 FK 公式假定的“零位”构型对应的角）。但真实电机的编码器读数 $\theta_{raw}$ 有自己的\emph{电机零点}（出厂 / 标定决定），二者一般\emph{不重合}。连接实物与模型靠一个标定偏置\cite{yu2021quadruped}：
\[
  \theta_{model}=\theta_{raw}+\theta_{offset}.
\]
工程要点（凝练不展开）：
\begin{itemize}
  \item 上电后必须先做\emph{关节零位标定}（如靠机械限位 / 外部工装确定 $\theta_{offset}$），否则 FK 算出的足端位置整体偏移，控制器全盘失效；
  \item 标定误差直接表现为站立姿态歪斜、足端落点系统性偏移——排查机器人“站不正”时优先检查 $\theta_{offset}$；
  \item 左右腿、不同电机的 $\theta_{offset}$ 各异，须逐关节标定并持久化存储。
\end{itemize}
这属于把抽象运动学落到真机的“最后一公里”，具体框架实现见 \cref{ch:quad_prac}。
\end{practice}

\begin{insight}[K13：零位 / 轴向耦合进几何宏时不可复用（电机零点$\neq$模型零点的工程放大）]\label{ins:qk-geom-macro-reuse}
\cref{prac:qk-motor-zero} 的“电机零点 $\neq$ 模型零点”是\emph{一个关节}的标定问题；把同一道理放到\emph{整机建模文件}的层面，就得到移植开源框架时一类极隐蔽的陷阱：\textbf{当一个几何宏把某条具体的“零位约定”硬编码进关节的安装位姿（origin）时，它就不再是机器人无关的通用件，对零位约定不同的机器人不可复用}。

以一套常见的四足建模描述为例。它把单腿写成一个可参数化的腿宏，整机靠该宏实例化四次（左右、前后各一份，详见 \cref{ch:quad_advanced} 的移植清单）。这个腿宏在书写膝关节（knee，记 KFE）相对大腿（thigh）的安装位姿时，直接令
\begin{equation}
  {}^{\text{thigh}}\boldsymbol{p}_{\text{KFE}}=[\,0,\ 0,\ -l_{\text{thigh}}\,]^\top,
  \label{eq:qk-kfe-origin}
\end{equation}
即代码里的 \texttt{KFE origin xyz="0 0 -thigh\_length"}。这条 origin 之所以把膝轴放在大腿正下方（沿 $-z$ 偏移一个大腿长 $l_{\text{thigh}}$），\emph{隐含假设了该机的零位构型是“大腿竖直向下”}。轴向同理被一并写死：髋外摆轴取 $(1,0,0)$、髋俯仰与膝俯仰轴取 $(0,1,0)$，四腿全同，左右之分靠站姿角而非翻轴实现（翻轴会双重取反、反而破坏对称）。\cref{eq:qk-kfe-origin} 与这套轴向\emph{只对“竖直向下零位”的机器人成立}。

问题在于：零位约定本是\emph{自定义量}。若目标机器人采用的是“大腿水平向后”零位（即与上例差一个绕 $+y$ 轴的 $+\pi/2$，\cref{tab:qk-zero-conv}），却照搬该腿宏，则 \cref{eq:qk-kfe-origin} 把膝轴放到了错误的方向上——整条腿的关节系朝向、乃至挂在其上的 STL 网格朝向\emph{系统性地转了 $90^\circ$}，FK 算出的足端位置全盘错位。这与 \cref{ins:qp-collision}（碰撞网格用错）是同源的失效：都是“把一份本该按本机几何重新生成的东西，当成通用件直接搬来”。

\begin{table}[H]\centering\small
\caption{两种大腿零位约定（同一描述的 origin 不能跨约定复用）}
\label{tab:qk-zero-conv}
\begin{tabular}{p{0.20\textwidth} p{0.30\textwidth} p{0.40\textwidth}}
\toprule
零位构型 & 大腿指向（机身系） & 复用 \cref{eq:qk-kfe-origin} 的后果 \\
\midrule
竖直向下 & $-z$（沿重力） & origin 自洽，FK 正确 \\
水平向后 & $-x$（与上者差绕 $+y$ 轴 $+\pi/2$） & 膝轴朝向错 $90^\circ$，STL 网格与足端位置全盘转向 \\
\bottomrule
\end{tabular}
\end{table}

\textbf{判定准则}：移植同作者的参考件时，与机器人无关的薄包（如 Gazebo 插件、IMU、传动、材质宏）可放心复用；但\emph{凡与零位 / 轴向约定耦合的几何宏（典型即腿的运动学宏）一律不可复用，必须按本机几何重新建}，其关节零位优先由 FK 实算给出（如先验地令足端落在髋正下方、四足质心对称且共面），而非沿用参考件的解析零位值。这条陷阱藏在关节 origin 一行里，仅核对“连杆 / 关节命名是否对齐参考”时极易漏过；完整的移植清单与“什么绝不能复用”的判据见 \cref{ch:quad_advanced}。
\end{insight}

% ════════════════════════════════════════════════════════════════
\section{常见误解辨析}
\label{sec:qk-misconceptions}

\begin{table}[H]\centering\small
\begin{tabular}{p{0.46\textwidth} p{0.46\textwidth}}
\toprule
\textbf{常见误解} & \textbf{正确理解} \\
\midrule
公式比坐标系重要 & 有运动学必先建坐标系；上下标只在回答“哪个系、从哪看”（\cref{ins:qk-frames-first}）\\
旋转矩阵就一个用途 & 三大用途：描述姿态 / 坐标变换 / 向量旋转（\cref{ins:qk-three-uses}）\\
旋转矩阵求逆用 \texttt{inverse()} & 逆 $=$ 转置，用 \texttt{transpose()}（\cref{eq:qk-R-orth}）\\
欧拉角适合长期描述姿态 & 有万向锁、不可插值；只配人机交互（\cref{ins:qk-euler-only-hmi}）\\
FK 和 IK 一样难 & FK 唯一解、IK 多解 / 无解 / 奇异（\cref{ins:qk-fk-easy-ik-hard}）\\
足端与关节角关系可线性化为 $\mathbf{A}\boldsymbol{q}$ & 位置非线性；唯有\emph{速度}层 $\dot{\boldsymbol{p}}=\mathbf{J}\dot{\boldsymbol{q}}$ 线性（\cref{ins:qk-diff-linearize}）\\
速度与力各用各的矩阵 & 同一 $\mathbf{J}$ 经功率守恒贯通：$\boldsymbol{\tau}=\mathbf{J}^\top\mathbf{F}$（\cref{thm:qk-duality}）\\
浮动基座要特殊处理 & 当成虚拟 $6$-DoF 关节 $\mathbf{S}_1=\mathbf{I}_6$，统一进树（\cref{ins:qk-floating-virtual}）\\
零空间运动会扰动主任务 & 零空间内运动末端不动，是“免费”内部运动（\cref{ins:qk-nullspace-free}）\\
任意右逆都行 & MP 伪逆给最小范数（能量最优）唯一解（\cref{thm:qk-mp-minnorm}）\\
$\mathbf{J}\mathbf{J}^\top$ 可逆就放心求逆 & 近奇异病态须阻尼最小二乘（\cref{ins:qk-damped}）\\
机身系随机身运动 & 取瞬时重合的静止惯性系，才能分解牵连 $+$ 关节速度（\cref{ins:qk-static-body}）\\
编码器读数就是模型关节角 & 差一个标定偏置 $\theta_{model}=\theta_{raw}+\theta_{offset}$（\cref{prac:qk-motor-zero}）\\
\bottomrule
\end{tabular}
\end{table}

% ════════════════════════════════════════════════════════════════
\section*{本章小结}

\begin{table}[H]\centering\small
\caption{符号表}
\begin{tabular}{lll}
\toprule
符号 & 含义 & 首次出现 \\
\midrule
${}^{\mathcal{A}}\mathbf{R}_{\mathcal{B}}\in\SO(3)$ & $\mathcal{A}$ 系看 $\mathcal{B}$ 系的旋转矩阵 & \cref{eq:qk-R-def} \\
$\boldsymbol{\Theta}=[\phi,\theta,\psi]^\top$ & ZYX 欧拉角（roll/pitch/yaw） & \cref{eq:qk-euler-to-R} \\
$(\cdot)^\wedge,(\cdot)^\vee$ & 反对称化 / 其逆 & \cref{eq:qk-skew} \\
$\boldsymbol{\omega}_w,\boldsymbol{\omega}_b$ & 世界系 / 机体系角速度 & \cref{note:qk-omega-frame} \\
$\mathbf{T}\in\SE(3)$ & 齐次变换矩阵 & \cref{eq:qk-T} \\
$\boldsymbol{q}_1,\mathbf{S}_1$ & 浮动基座虚拟关节变量 / 子空间 $\mathbf{I}_6$ & \cref{eq:qk-floating-joint} \\
$\lambda(i)$ & 刚体 $i$ 的父节点 & \cref{eq:qk-lambda} \\
$\theta_1,\theta_2,\theta_3$ & 单腿三关节角（abad/hip/knee） & \cref{eq:qk-3d-T} \\
$l_1,l_2,l_3$ & abad 偏置 / 大腿 / 小腿长 & \cref{exam:qk-dims} \\
$\mathbf{J}\in\mathbb{R}^{3\times3}$ & 单腿足端雅可比 & \cref{eq:qk-jacobian-def} \\
$\boldsymbol{\tau},\mathbf{F}$ & 关节力矩 / 足端力 & \cref{thm:qk-duality} \\
$\boldsymbol{v},\hat{\boldsymbol{f}}$ & 六维空间运动 / 力向量 & \cref{eq:qk-spatial-vel,eq:qk-spatial-force} \\
$\mathbf{X},\mathbf{X}^{*}$ & Plücker 运动 / 力坐标变换 & \cref{eq:qk-plucker-X} \\
$\boldsymbol{v}\times,\boldsymbol{v}\times^{*}$ & 空间叉乘（运动 / 力） & \cref{eq:qk-spatial-cross} \\
$\hat{\mathbf{I}}$ & 空间惯量 & \cref{eq:qk-spatial-inertia} \\
${}_{b}\mathbf{J}_i,{}_{b}^{\mathcal{Q}}\mathbf{J}_i$ & 刚体 / 任务雅可比 & \cref{eq:qk-rigid-jacobian,eq:qk-task-jacobian} \\
$\mathcal{N}(\mathbf{J})$ & 雅可比零空间 & \cref{eq:qk-null-def} \\
$\mathbf{J}^{+},\mathbf{N}$ & MP 伪逆 / 零空间投影 & \cref{eq:qk-general-solution} \\
$\lambda$（阻尼） & 阻尼最小二乘因子 & \cref{eq:qk-damped-pinv} \\
$\boldsymbol{v}_{sfB}$ & 足端相对机身速度（可测量） & \cref{eq:qk-foot-vel-world} \\
$\theta_{offset}$ & 电机零点标定偏置 & \cref{prac:qk-motor-zero} \\
\bottomrule
\end{tabular}
\end{table}

\begin{table}[H]\centering\small
\caption{公式速查表}
\begin{tabular}{lll}
\toprule
公式 / 结论 & 一句话说明 & 对应 \\
\midrule
旋转矩阵定义 \cref{eq:qk-R-def} & 列即 $\mathcal{B}$ 轴在 $\mathcal{A}$ 系坐标 & 下标消去 \cref{eq:qk-R-cancel} \\
角速度张量 \cref{eq:qk-Rdot-world} & $\dot{\mathbf{R}}\mathbf{R}^\top=\boldsymbol{\omega}^\wedge$（反对称） & \cref{note:qk-omega-frame} \\
欧拉角逆解 \cref{eq:qk-R-to-euler} & atan2 三式，俯仰奇异 & \cref{ins:qk-euler-only-hmi} \\
浮动基座关节 \cref{eq:qk-floating-joint} & $\mathbf{S}_1=\mathbf{I}_6$ & \cref{ins:qk-floating-virtual} \\
单腿 FK \cref{eq:qk-3d-fk} & 齐次变换连乘闭式 & \cref{exam:qk-fk-num} \\
单腿 IK 三步 \cref{eq:qk-ik-theta1,eq:qk-ik-theta3,eq:qk-ik-theta2} & $\theta_1$ 投影 $/\theta_3$ 余弦 $/\theta_2$ 整理 & \cref{exam:qk-ik-num} \\
足端雅可比 \cref{eq:qk-jacobian-def} & $\dot{\boldsymbol{p}}=\mathbf{J}\dot{\boldsymbol{q}}$（一阶微分变线性） & \cref{ins:qk-diff-linearize} \\
静力学对偶 \cref{thm:qk-duality} & $\boldsymbol{\tau}=\mathbf{J}^\top\mathbf{F}$（功率守恒） & \cref{ins:qk-duality} \\
Plücker 变换 \cref{eq:qk-plucker-X-block} & $\mathbf{X}^{*}=\mathbf{X}^{-\top}$（保功率） & 运动 / 力对偶 \\
空间运动方程 \cref{eq:qk-spatial-eom} & $\hat{\mathbf{f}}=\hat{\mathbf{I}}\boldsymbol{a}+\boldsymbol{v}\times^{*}\hat{\mathbf{I}}\boldsymbol{v}$ & \cref{ins:qk-eom-contrast} \\
递推 FK \cref{eq:qk-recursive} & $\boldsymbol{v}_i=\boldsymbol{v}_{\lambda(i)}+\mathbf{S}_i\dot{\boldsymbol{q}}_i$ & \cref{alg:qk-recursive-fk} \\
通解结构 \cref{eq:qk-general-solution} & 特解 $+$ 零空间内部运动 & \cref{ins:qk-nullspace-free} \\
MP 伪逆 \cref{eq:qk-mp-result} & $\mathbf{J}^\top(\mathbf{J}\mathbf{J}^\top)^{-1}$ 最小范数 & \cref{thm:qk-mp-minnorm} \\
阻尼伪逆 \cref{eq:qk-damped-pinv} & $+\lambda^2\mathbf{I}$ 近奇异稳健 & \cref{ins:qk-damped} \\
整机足端速度 \cref{eq:qk-foot-vel-world} & 牵连 $+$ 关节，$\boldsymbol{v}_{sfB}$ 可测 & \cref{ins:qk-static-body} \\
\bottomrule
\end{tabular}
\end{table}

\begin{table}[H]\centering\small
\caption{知识点总表}
\begin{tabular}{clll}
\toprule
\# & 知识点 & 核心要点 & 难度 \\
\midrule
1 & 坐标系体系 & 有运动学必先建系；Body/Joint/Link 分层 & \(\bigstar\bigstar\) \\
2 & 旋转矩阵三用途 & 描述姿态 / 坐标变换 / 向量旋转 & \(\bigstar\bigstar\) \\
3 & 欧拉角与逆解 & ZYX、atan2、万向锁、只配人机交互 & \(\bigstar\bigstar\) \\
4 & 角速度张量 & $\dot{\mathbf{R}}\mathbf{R}^\top=\boldsymbol{\omega}^\wedge$ 反对称严格推 & \(\bigstar\bigstar\bigstar\) \\
5 & 浮动基座虚拟关节 & $\mathbf{S}_1=\mathbf{I}_6$，唯一结构差别 & \(\bigstar\bigstar\bigstar\) \\
6 & 运动学树 & 13 刚体、父节点集、统一递推 & \(\bigstar\bigstar\bigstar\) \\
7 & 单腿 FK & 齐次变换连乘闭式 & \(\bigstar\bigstar\) \\
8 & 单腿 IK 解耦三步 & $\theta_1$ 投影 / $\theta_3$ 余弦 / $\theta_2$ 整理 & \(\bigstar\bigstar\bigstar\) \\
9 & 足端雅可比 & 一阶微分把非线性变线性 & \(\bigstar\bigstar\bigstar\) \\
10 & 静力学对偶 & 功率守恒 $\Rightarrow\boldsymbol{\tau}=\mathbf{J}^\top\mathbf{F}$ & \(\bigstar\bigstar\bigstar\) \\
11 & Plücker 空间向量 & 运动 / 力六维、$\mathbf{X}/\mathbf{X}^{*}$、空间惯量 & \(\bigstar\bigstar\bigstar\bigstar\) \\
12 & 递推 FK / 全身雅可比 & $O(N)$ 递推、$\mathbf{J}_1\dots\mathbf{J}_4$ & \(\bigstar\bigstar\bigstar\bigstar\) \\
13 & 零空间与投影 & 免费内部运动、$\mathbf{N}=\mathbf{I}-\mathbf{J}^{+}\mathbf{J}$ & \(\bigstar\bigstar\bigstar\bigstar\) \\
14 & MP 伪逆证明 & 拉格朗日法、最小范数 & \(\bigstar\bigstar\bigstar\bigstar\) \\
15 & 奇异与阻尼 & 病态须阻尼最小二乘 & \(\bigstar\bigstar\bigstar\) \\
16 & 整机运动学 & 姿态$\leftrightarrow$足端、牵连 $+$ 关节速度 & \(\bigstar\bigstar\bigstar\) \\
\bottomrule
\end{tabular}
\end{table}

% ════════════════════════════════════════════════════════════════
\section*{延伸阅读}
\begin{itemize}
  \item \textbf{一手详尽推导}：于宪元《基于稳定性的仿生四足机器人控制系统设计》\cite{yu2021quadruped}（北航硕论，2021）——§2.2 坐标系与角速度张量、第五章 Plücker 空间向量与运动学树 / 全身雅可比，本章理论主体据此展开。
  \item \textbf{工程落地与单腿运动学}：宇树《四足机器人控制算法》\cite{unitree2023quadruped}——旋转矩阵下标消去、单腿 FK/IK 解耦三步、足端雅可比与静力学对偶、整机姿态$\leftrightarrow$足端，配套 \texttt{unitree\_guide} 开源代码。
  \item \textbf{多刚体动力学统一框架}：Featherstone\cite{featherstone2008rigid}——空间向量代数（Plücker 坐标、空间叉乘、空间惯量、递推算法）的奠基专著，本章 \cref{sec:qk-plucker} 的理论源头。
  \item \textbf{现代机器人学几何视角}：Lynch \& Park\cite{lynch2017modern}——旋量 / 指数坐标 / 伴随表示，与本章齐次变换、Plücker 变换可对照阅读。
  \item \textbf{李群几何地基}：见 \cref{ch:lie}（$\SO(3)/\SE(3)$、指数 / 对数映射、右扰动）、\cref{ch:rigid_body}（刚体运动与欧拉角）、\cref{ch:matderiv}（雅可比 / 矩阵求导）。
\end{itemize}

\section*{本章与后续章节的关系}
\begin{table}[H]\centering\small
\begin{tabular}{lll}
\toprule
相关章节 & 关系 & 本章接口 \\
\midrule
\cref{ch:lie} 李群与李代数 & 提供 $\dot{\mathbf{R}}=\mathbf{R}\boldsymbol{\omega}_b^\wedge$、$(\cdot)^\wedge$、指数坐标 & \cref{eq:qk-Rdot-world,eq:qk-rodrigues} \\
\cref{ch:rigid_body} 刚体运动 & 提供 $\SE(3)$、ZYX 欧拉角 & \cref{eq:qk-T,eq:qk-euler-to-R} \\
\cref{ch:matderiv} 矩阵求导 & 提供雅可比、秩--零度定理 & \cref{eq:qk-jacobian-def,eq:qk-rank-nullity} \\
\cref{ch:nlopt} 优化 / QP & 提供拉格朗日法、正则最小二乘 & \cref{thm:qk-mp-minnorm,eq:qk-damped-pinv} \\
\cref{ch:quad_mpc} 单刚体 MPC & 消费 $\boldsymbol{\tau}=\mathbf{J}^\top\mathbf{f}$、惯量、陀螺项对照 & \cref{thm:qk-duality,ins:qk-eom-contrast} \\
\cref{ch:quad_wbc} 全身控制 & 消费零空间投影 / 伪逆 / 全身雅可比 & \cref{thm:qk-nullspace-proj,eq:qk-task-jacobian} \\
\cref{ch:quad_est} 足式状态估计 & 消费整机足端速度（前向模型） & \cref{ins:qk-vsfB} \\
\cref{ch:quad_prac} 实践 & 消费电机零点标定 & \cref{prac:qk-motor-zero} \\
\bottomrule
\end{tabular}
\end{table}

\section*{故障排查手册}
\begin{enumerate}
  \item \textbf{症状}：足端位置整体偏移 / 机身站不正。\textbf{原因}：关节零位标定不当，$\theta_{offset}$ 错（\cref{prac:qk-motor-zero}）。\textbf{对策}：重做逐关节零位标定，核对左右腿 abad 偏置符号（\cref{note:qk-ocr-leg-sign}）。
  \item \textbf{症状}：IK 返回 NaN / 关节角跳变。\textbf{原因}：目标点超出可达工作空间，$\arccos$ 参数越界（\cref{pit:qk-ik-multi}）；或选了超限位的另一支多解。\textbf{对策}：对 $\arccos$ 参数裁剪到 $[-1,1]$，做可达性预判，按限位选解支。
  \item \textbf{症状}：近奇异时关节速度爆炸 / 机器人抽搐。\textbf{原因}：$\mathbf{J}\mathbf{J}^\top$ 病态，直接求逆放大误差（\cref{ins:qk-damped}）。\textbf{对策}：改用阻尼最小二乘 \cref{eq:qk-damped-pinv}，按操作度 $w$ 调 $\lambda$。
  \item \textbf{症状}：姿态在大俯仰下乱跳 / 丢自由度。\textbf{原因}：欧拉角万向锁（$\theta\to\pm90^\circ$，\cref{ins:qk-euler-only-hmi}）。\textbf{对策}：姿态长期表达改用四元数 / 旋转矩阵，欧拉角仅用于输入 / 显示。
  \item \textbf{症状}：力矩方向错 / 机器人“顶”地或“瘫”地。\textbf{原因}：$\boldsymbol{\tau}=\mathbf{J}^\top\mathbf{F}$ 的力方向约定混淆（对地力 vs 地对足反力，\cref{thm:qk-duality} 后）。\textbf{对策}：核对符号，地对足反力须取 $\boldsymbol{\tau}=-\mathbf{J}^\top\mathbf{F}'$。
  \item \textbf{症状}：状态估计速度发散。\textbf{原因}：误用机身绝对速度 $\boldsymbol{v}_s$（不可测）而非足端相对机身速度 $\boldsymbol{v}_{sfB}$（\cref{ins:qk-vsfB}）。\textbf{对策}：以 $\boldsymbol{v}_{sfB}$ 为观测，由支撑足不滑移反推 $\boldsymbol{v}_s$，详见 \cref{ch:quad_est}。
\end{enumerate}

% ════════════════════════════════════════════════════════════════
\section{本章推导附录}
\label{sec:qk-appendix}

\begin{derivation}[单腿雅可比第二、三行的逐项求导]\label{deriv:qk-jacobian}
对 \cref{eq:qk-3d-fk} 的 $y_P,z_P$ 逐元求偏导，给出 $\mathbf{J}$ 第二、三行的完整推导。记 $s_{23}=\sin(\theta_2+\theta_3),c_{23}=\cos(\theta_2+\theta_3)$。
对 $y_P=-l_3 s_1 c_{23}+l_1 c_1-l_2 c_2 s_1$：
\[
  J_{21}=\frac{\partial y_P}{\partial\theta_1}=-l_3 c_1 c_{23}-l_1 s_1-l_2 c_1 c_2,\quad
  J_{22}=\frac{\partial y_P}{\partial\theta_2}=l_3 s_1 s_{23}+l_2 s_1 s_2,\quad
  J_{23}=\frac{\partial y_P}{\partial\theta_3}=l_3 s_1 s_{23}.
\]
对 $z_P=l_3 c_1 c_{23}+l_1 s_1+l_2 c_1 c_2$：
\[
  J_{31}=\frac{\partial z_P}{\partial\theta_1}=-l_3 s_1 c_{23}+l_1 c_1-l_2 s_1 c_2,\quad
  J_{32}=\frac{\partial z_P}{\partial\theta_2}=-l_3 c_1 s_{23}-l_2 c_1 s_2,\quad
  J_{33}=\frac{\partial z_P}{\partial\theta_3}=-l_3 c_1 s_{23}.
\]
连同 \cref{eq:qk-jacobian-row1} 的第一行 $(0,\ l_3 c_{23}+l_2 c_2,\ l_3 c_{23})$，即得完整 $3\times3$ 足端雅可比。读者可在 \cref{exam:qk-fk-num} 的构型 $(\theta_1,\theta_2,\theta_3)=(0,-\tfrac{\pi}{4},\tfrac{\pi}{2})$ 代入数值，验证 $\det\mathbf{J}\neq0$（远奇异，可直接求逆）。
\end{derivation}

\begin{derivation}[Plücker 力变换 $\mathbf{X}^{*}=\mathbf{X}^{-\top}$ 由保功率导出（\cref{eq:qk-plucker-X}）]\label{deriv:qk-plucker-dual}
功率是运动向量与力向量的内积 $P=\boldsymbol{m}^\top\hat{\boldsymbol{f}}$，它是\emph{标量}，与坐标系无关，故 ${}^{\mathcal{A}}\boldsymbol{m}^\top\,{}^{\mathcal{A}}\hat{\boldsymbol{f}}={}^{\mathcal{B}}\boldsymbol{m}^\top\,{}^{\mathcal{B}}\hat{\boldsymbol{f}}$。代入运动变换 ${}^{\mathcal{B}}\boldsymbol{m}={}^{\mathcal{B}}\mathbf{X}_{\mathcal{A}}\,{}^{\mathcal{A}}\boldsymbol{m}$ 与待定的力变换 ${}^{\mathcal{B}}\hat{\boldsymbol{f}}=\mathbf{Y}\,{}^{\mathcal{A}}\hat{\boldsymbol{f}}$：
\[
  {}^{\mathcal{A}}\boldsymbol{m}^\top\,{}^{\mathcal{A}}\hat{\boldsymbol{f}}
  =({}^{\mathcal{B}}\mathbf{X}_{\mathcal{A}}\,{}^{\mathcal{A}}\boldsymbol{m})^\top(\mathbf{Y}\,{}^{\mathcal{A}}\hat{\boldsymbol{f}})
  ={}^{\mathcal{A}}\boldsymbol{m}^\top\,({}^{\mathcal{B}}\mathbf{X}_{\mathcal{A}}^\top\mathbf{Y})\,{}^{\mathcal{A}}\hat{\boldsymbol{f}}.
\]
对任意 ${}^{\mathcal{A}}\boldsymbol{m},{}^{\mathcal{A}}\hat{\boldsymbol{f}}$ 成立要求 ${}^{\mathcal{B}}\mathbf{X}_{\mathcal{A}}^\top\mathbf{Y}=\mathbf{I}$，即 $\mathbf{Y}={}^{\mathcal{B}}\mathbf{X}_{\mathcal{A}}^{-\top}={}^{\mathcal{B}}\mathbf{X}_{\mathcal{A}}^{*}$。这就是力变换为运动变换的逆转置的来由——运动与力的对偶性在能量层面的体现，与 \cref{thm:qk-duality} 的功率守恒同源。
\end{derivation}

\begin{derivation}[阻尼伪逆等价于正则最小二乘（\cref{eq:qk-damped-pinv}）]\label{deriv:qk-damped}
考虑正则最小二乘 $\min_{\dot{\boldsymbol{q}}}\ \tfrac12\|\mathbf{J}\dot{\boldsymbol{q}}-\dot{\boldsymbol{x}}\|^2+\tfrac12\lambda^2\|\dot{\boldsymbol{q}}\|^2$（\cref{ch:nlopt}）。对 $\dot{\boldsymbol{q}}$ 求梯度置零：
\[
  \mathbf{J}^\top(\mathbf{J}\dot{\boldsymbol{q}}-\dot{\boldsymbol{x}})+\lambda^2\dot{\boldsymbol{q}}=\mathbf{0}
  \ \Rightarrow\ (\mathbf{J}^\top\mathbf{J}+\lambda^2\mathbf{I})\dot{\boldsymbol{q}}=\mathbf{J}^\top\dot{\boldsymbol{x}}.
\]
由矩阵恒等式 $(\mathbf{J}^\top\mathbf{J}+\lambda^2\mathbf{I})^{-1}\mathbf{J}^\top=\mathbf{J}^\top(\mathbf{J}\mathbf{J}^\top+\lambda^2\mathbf{I})^{-1}$（推挽 / push-through 恒等式），得
\[
  \dot{\boldsymbol{q}}=\mathbf{J}^\top(\mathbf{J}\mathbf{J}^\top+\lambda^2\mathbf{I})^{-1}\dot{\boldsymbol{x}},
\]
即 \cref{eq:qk-damped-pinv} 的阻尼伪逆。当 $\lambda\to0$ 退回 MP 伪逆（\cref{eq:qk-mp-result}）；$\lambda>0$ 时 $\mathbf{J}\mathbf{J}^\top+\lambda^2\mathbf{I}$ 特征值有下界 $\lambda^2$ 故始终可逆，这就是近奇异稳健化的数学本质——用一点跟踪误差换条件数的有界性。
\end{derivation}
